\chapter{完全信息动态博弈}\label{chap:02}

% ============================================================================
%  2.1  Rotten Kid 定理（Becker, 1974）
% ============================================================================
\exer{2.1}
本题取自首先由贝克尔（Becker，1974）提出并分析的一个模型。假设一个家长和他的孩子进行如下的博弈：第一，小孩选择一个行动 $A$，可使孩子获得收入 $I_c(A)$，并使家长得到收入 $I_p(A)$（可以认为 $I_c(A)$ 为孩子减去因 $A$ 发生的各种成本后的净收益）；第二，家长观测到收入 $I_c$ 和 $I_p$，然后选择给孩子的奖励或惩罚 $B$。孩子的收益为 $U(I_c+B)$，家长的收益为 $V(I_p-B)+k(I_c+B)$，其中 $k>0$ 反映出家长关心孩子的福利。假定行动是一个非负数字 $A\ge 0$，收入函数 $I_c(A)$ 和 $I_p(A)$ 为严格凹且分别在 $A_c>0$ 和 $A_p>0$ 达到最大值；奖励或惩罚 $B$ 为或正或负的数字；且效用函数 $U$ 和 $V$ 递增并严格凹。证明“宠坏的孩子”（Rotten Kid）定理：在逆向归纳解中，孩子选择可使全家收入 $I_c(A)+I_p(A)$ 最大化的行为，尽管在效用函数中，只有家长显示出利他的特点。

\sol
博弈时序为两阶段：第 1 阶段孩子选择 $A\ge 0$；第 2 阶段家长观测到 $(I_c(A),I_p(A))$ 后选择转移 $B\in\mathbb{R}$（$B>0$ 为奖励，$B<0$ 为惩罚）。

\textbf{第 2 阶段（家长的最优转移）}。给定 $A$，家长选择 $B$ 最大化
\[
W(B)=V\!\bigl(I_p(A)-B\bigr)+k\bigl(I_c(A)+B\bigr).
\]\indent
由于 $U,V$ 递增且严格凹，$V$ 严格递增意味着 $V'$ 存在且为正、严格递减。对 $B$ 求导：
\[
\frac{dW}{dB}=-V'\!\bigl(I_p(A)-B\bigr)+k.
\]\indent
令其为零得家长最优转移 $B^*(A)$ 满足
\[
V'\!\bigl(I_p(A)-B^*(A)\bigr)=k.
\]\indent
记方程 $V'(x)=k$ 的唯一解为 $\bar x$（因为 $V'$ 严格递减、值域覆盖 $k$），则家长把自身的消费维持在 $\bar x$：
\[
I_p(A)-B^*(A)=\bar x \quad\Longrightarrow\quad B^*(A)=I_p(A)-\bar x.
\]\indent
注意 $\bar x$ 只取决于 $k$ 与 $V$，与孩子的行动 $A$ 无关。

\textbf{第 1 阶段（孩子的最优化）}。孩子预期到家长的转移规则 $B^*(A)=I_p(A)-\bar x$，其第 2 阶段收益为
\[
U\!\bigl(I_c(A)+B^*(A)\bigr)=U\!\bigl(I_c(A)+I_p(A)-\bar x\bigr).
\]\indent
由于 $U$ 严格递增，孩子选择 $A$ 等价于最大化
\[
I_c(A)+I_p(A)-\bar x,
\]\indent
即最大化全家总收入 $I_c(A)+I_p(A)$。而 $I_c+I_p$ 为严格凹函数（两严格凹函数之和仍严格凹），故存在唯一最大值点 $A^*$，满足
\[
I_c'(A^*)+I_p'(A^*)=0.
\]

\textbf{结论}。尽管孩子完全自私（其效用只取决于自身收入），家长出于利他动机的转移 $B^*(A)$ 使孩子意识到：自己的行为会通过影响 $I_p(A)$ 而改变所收到的转移，从而孩子的最优选择恰好是使全家总收入最大化的 $A^*$。这正是“宠坏的孩子”定理：逆向归纳解中孩子选择 $A^*$ 使 $I_c+I_p$ 最大化。\cmt{关键在于家长的最优转移使“孩子的边际收益”等于全家收入的边际收益，自私的孩子被“收买”为家庭收入最大化者。}

% ============================================================================
%  2.2  Samaritan's Dilemma（Buchanan, 1975）
% ============================================================================
\exer{2.2}
现在假设家长和孩子进行一个不同的博弈，首先由布坎南（Buchanan，1975）所分析。令收入 $I_c$ 和 $I_p$ 为外生给定的，第一，孩子决定收入 $I_c$ 中的多少用于储蓄（$S$）以备将来，并消费掉其余部分 $I_c-S$。第二，家长观测到孩子的选择 $S$ 并决定一个赠与额 $B$。孩子的收益为当期和未来的效用之和 $U_1(I_c-S)+U_2(S+B)$；家长的收益为 $V(I_p-B)+[U_1(I_c-S)+U_2(S+B)]$。假定效用函数 $U_1$、$U_2$ 和 $V$ 递增并且严格凹，证明存在“乐善好施俘论”（Samaritan's Dilemma）：在逆向归纳解中，孩子的储蓄非常少，从而可诱使家长给予更高的赠与（即如果 $S$ 增加，并使 $B$ 相应减少，家长和孩子的福利都会提高）。

\sol
博弈时序：第 1 阶段孩子选择储蓄 $S\ge 0$；第 2 阶段家长观测到 $S$ 后选择赠与 $B$。

\textbf{第 2 阶段（家长对给定的 $S$ 选择 $B$）}。家长最大化
\[
\Pi_p(B;S)=V(I_p-B)+U_2(S+B)+U_1(I_c-S),
\]\indent
其中 $U_1(I_c-S)$ 与 $B$ 无关。一阶条件：
\[
\frac{\partial\Pi_p}{\partial B}=-V'(I_p-B)+U_2'(S+B)=0.
\]\indent
记由此决定的最优赠与为 $B^*(S)$，它满足
\[
V'\!\bigl(I_p-B^*(S)\bigr)=U_2'\!\bigl(S+B^*(S)\bigr).
\]\indent
对隐函数 $F(B,S)\equiv V'(I_p-B)-U_2'(S+B)=0$ 求偏导：
\[
\frac{\partial F}{\partial B}=-V''(I_p-B)-U_2''(S+B)>0,\qquad
\frac{\partial F}{\partial S}=-U_2''(S+B)>0.
\]\indent
故
\[
\frac{dB^*(S)}{dS}=-\frac{\partial F/\partial S}{\partial F/\partial B}<0.
\]\indent
即孩子的储蓄越高，家长的最优赠与越低——这是“乐善好施”家长的理性反应。

\textbf{第 1 阶段（孩子选择 $S$）}。孩子预期 $B^*(S)$，最大化
\[
\Pi_k(S)=U_1(I_c-S)+U_2\!\bigl(S+B^*(S)\bigr).
\]\indent
一阶条件：
\[
\frac{d\Pi_k}{dS}=-U_1'(I_c-S)+U_2'\!\bigl(S+B^*(S)\bigr)
+U_2''\!\bigl(S+B^*(S)\bigr)\frac{dB^*(S)}{dS}=0.
\]\indent
由于 $U_2''<0$ 且 $dB^*/dS<0$，末项 $U_2''\cdot dB^*/dS>0$，于是
\[
-U_1'(I_c-S)+U_2'\!\bigl(S+B^*(S)\bigr)=-\underbrace{U_2''\frac{dB^*}{dS}}_{>0}<0,
\]\indent
即
\[
U_2'\!\bigl(S+B^*(S)\bigr)<U_1'(I_c-S).\tag{$\star$}
\]

\textbf{与帕累托有效解的对比}。若家长与孩子作为一个整体选择 $(S,B)$ 以最大化联合福利
\[
J(S,B)=U_1(I_c-S)+U_2(S+B)+V(I_p-B),
\]\indent
在给定 $S$ 下对 $B$ 的最优条件恰与家长相同（即 $B^*(S)$ 已是最优），但对 $S$ 的最优条件为
\[
\frac{\partial J}{\partial S}=-U_1'(I_c-S)+U_2'\!\bigl(S+B^*(S)\bigr)=0.
\]\indent
而孩子在 $(\star)$ 中多出一个正项，因此孩子的均衡储蓄 $S^{BI}$ 满足 $U_2'(S^{BI}+B^*)<U_1'(I_c-S^{BI})$，即相对有效水平“储蓄不足”。换言之，孩子知道提高 $S$ 会诱使家长减少赠与 $B$（因 $dB^*/dS<0$），于是故意少储蓄以套取更多赠与。由于 $B^*(S)$ 随 $S$ 下降，少储蓄反而对应更高的 $B$。若外生地提高 $S$ 并相应调低 $B$（保持家长福利不变），则 $(U_1+U_2)$ 会上升——即存在帕累托改进。这便是“乐善好施俘论”：家长的慷慨承诺反而使孩子减少储蓄，双方本可更好却陷入低效均衡。

% ============================================================================
%  2.3  不同贴现因子的鲁宾斯坦讨价还价
% ============================================================================
\exer{2.3}
假定鲁宾斯坦的无期限讨价还价博弈中，参与者的贴现因子不同，参与者 1 为 $\delta_1$，参与者 2 为 $\delta_2$。运用本书中的论证方法证明在逆向归纳解中，参与者 1 向 2 提出的解决方案为
\[
\bigl(x_1^*,\,1-x_1^*\bigr),\qquad
x_1^*=\frac{1-\delta_2}{1-\delta_1\delta_2},
\]\indent
并被参与者 2 接受（即参与者 1 得到 $x_1^*$，参与者 2 得到 $1-x_1^*$）。

\sol
无限期鲁宾斯坦博弈：奇数期由参与者 1 出价 $(x,1-x)$（1 得 $x$，2 得 $1-x$），偶数期由参与者 2 出价；任一方均可接受或拒绝，拒绝则进入下一期。参与者 1、2 的贴现因子分别为 $\delta_1,\delta_2$。

设 $a$ 为“轮到参与者 1 出价时，参与者 1 的均衡收益”（即 1 作为提议方时的价值），$b$ 为“轮到参与者 2 出价时，参与者 1 的均衡收益”（即 1 作为回应方时的价值）。由无限期博弈的对称性，这些值在每一轮相同。

\textbf{参与者 2 出价时（偶数期）}。参与者 2 要使参与者 1 接受，必须让 1 当期所得不低于其拒绝后下一期（参与者 1 出价）所得，即
\[
1\text{ 的份额}\;\ge\;\delta_1\cdot a.
\]\indent
参与者 2 把 1 的份额恰好定为 $\delta_1 a$，自己保留 $1-\delta_1 a$。故当参与者 2 出价时，参与者 1 得到
\[
b=\delta_1 a.\tag{1}
\]

\textbf{参与者 1 出价时（奇数期）}。参与者 1 要使参与者 2 接受，须让 2 当期所得不低于其拒绝后下一期（参与者 2 出价）所得。参与者 2 出价时 2 自己的收益为 $1-b$（因为总额归一化为 1，且那时 1 得 $b$），贴现到本期为 $\delta_2(1-b)$。因此参与者 1 给 2 的份额至少应为 $\delta_2(1-b)$，于是
\[
1-x_1^*=\delta_2(1-b)\quad\Longrightarrow\quad
x_1^*=1-\delta_2(1-b)=1-\delta_2+\delta_2 b.
\]\indent
代入 $(1)$ 的 $b=\delta_1 a$：
\[
a=1-\delta_2+\delta_1\delta_2 a
\;\Longrightarrow\;
a(1-\delta_1\delta_2)=1-\delta_2
\;\Longrightarrow\;
a=x_1^*=\frac{1-\delta_2}{1-\delta_1\delta_2}.
\]\indent
从而参与者 2 所得为
\[
1-x_1^*=\frac{\delta_2(1-\delta_1)}{1-\delta_1\delta_2}.
\]

\textbf{验证对称性特例}。若 $\delta_1=\delta_2=\delta$，则
\[
x_1^*=\frac{1-\delta}{1-\delta^2}=\frac{1}{1+\delta},
\]\indent
即教材正文中的标准结果，参与者 1 出价 $(1/(1+\delta),\,\delta/(1+\delta))$ 并被接受。

\textbf{结论}。逆向归纳（子博弈精炼）解为：参与者 1 在第 1 期提出分配 $\bigl(x_1^*,1-x_1^*\bigr)$，参与者 2 接受；此后任何偏离都会触发对方在下一期索要同样苛刻的份额，故无一方愿偏离。$\blacksquare$

% ============================================================================
%  2.4  两阶段合伙人贡献博弈
% ============================================================================
\exer{2.4}
两合伙人希望能完成一个项目，在项目结束时，每一合伙人得到收益 $V$，但结束前则一无所得，尚需 $R$ 的成本。两合伙人都不能承诺只靠自己的力量完成项目，于是他们决定进行如下的两阶段博弈：在第一阶段，合伙人 1 选择为完成项目所作的贡献 $c_1$，如果这一贡献足以令项目完成，则博弈结束，每一参与者得到收益 $V$；如果这一贡献不足以令项目完成（即 $c_1<R$），则在第二阶段合伙人 2 选择为完成项目所作的贡献 $c_2$，如果（不考虑贴现的）两个贡献之和足以完成项目，则博弈结束，每一参与者得到 $V$；如果贡献之和不足以完成项目，则博弈结束，两参与者所得的收益均为 $0$。

每一合伙人必须从其他可带来收益的活动中抽取部分资金，投入到该项目，这样做的最优方法是从先收益最低的其他活动中抽资，结果使贡献的（机会）成本为贡献大小的凸函数。假设对每一参与者贡献 $c$ 的成本为 $c^2$，并假定参与者 1 对其第二阶段的收益用贴现因子 $\delta$ 进行贴现。针对三个参数 $\{V,R,\delta\}$ 的不同情况，分别计算出此两阶段贡献博弈惟一的逆向归纳解；参见阿德马蒂（Admati）和佩里（Perry）对无限重复情况的讨论。

\sol
参与人：合伙人 1、2。行动：第 1 阶段 $c_1\ge 0$；若 $c_1<R$，第 2 阶段 $c_2\ge 0$。收益规则：
\[
\begin{aligned}
\text{若 }c_1\ge R:\;& (V-c_1^2,\ V);\\
\text{若 }c_1<R,\ c_1+c_2\ge R:\;& \bigl(\delta V-c_1^2,\ V-c_2^2\bigr);\\
\text{若 }c_1<R,\ c_1+c_2<R:\;& \bigl(-c_1^2,\ -c_2^2\bigr).
\end{aligned}
\]\indent
（题设只让合伙人 1 对第二阶段收益贴现，故 $\delta$ 只作用于其第二阶段的 $V$；合伙人 2 不贴现。项目失败时两人均未得到 $V$，但仍各自承担已发生的机会成本 $c_i^2$。）
下面用逆向归纳法求解。

\textbf{第 2 阶段：给定 $c_1<R$，合伙人 2 的选择}。2 \emph{不}贴现，其利润为
\[
\pi_2=
\begin{cases}
V-c_2^2, & c_2\ge R-c_1,\\[4pt]
-c_2^2, & 0\le c_2<R-c_1.
\end{cases}
\]\indent
在“成功”区间 $c_2\ge R-c_1$，利润随 $c_2$ 增大而减小（成本上升），故 2 取最小可行值 $c_2=R-c_1$，得 $\pi_2=V-(R-c_1)^2$；在“失败”区间取 $c_2=0$ 得 $\pi_2=0$。因此
\[
c_2^*(c_1)=
\begin{cases}
R-c_1, & V>(R-c_1)^2,\\[3pt]
0\ \text{或}\ R-c_1, & V=(R-c_1)^2,\\[3pt]
0, & V<(R-c_1)^2.
\end{cases}
\tag{$\dagger$}
\]

\textbf{第 1 阶段：合伙人 1 的选择}。1 预见到 $(\dagger)$，比较三类可选贡献：

\emph{(i) 独自完成}：$c_1=R$，得 $\pi_1=V-R^2$。

\emph{(ii) 诱使 2 补足}：当 $0\le V\le R^2$ 时，临界贡献为 $c_1=R-\sqrt V$。在该点合伙人 2 对补足与不补足无差异；若其选择补足，合伙人 1 的收益为
\[
A(V)=\delta V-\bigl(R-\sqrt V\bigr)^2.
\]\indent
在严格诱使区间 $c_1>R-\sqrt V$，该收益随 $c_1$ 增加而下降。因而当均衡路径选择此方案时，必须由合伙人 2 在临界点选择补足；若不补足，且 $A(V)>0$，合伙人 1 的收益只有一个未达到的上确界，纯战略逆向归纳解不存在。

\emph{(iii) 无人补足}：$c_1=0$ 且 2 不贡献，得 $\pi_1=0$。

先看阈值 $\pi_1\ge0\iff\delta V\ge(R-\sqrt V)^2$。令 $s=\sqrt V\in(0,R]$，整理得
\[
\delta s^2-(R-s)^2=-(1-\delta)s^2+2Rs-R^2\ \ge 0
\iff \frac{R}{1+\sqrt\delta}\le s\le\frac{R}{1-\sqrt\delta}.
\]\indent
故
\[
\pi_1\ge 0\iff V\ge\frac{R^2}{(1+\sqrt\delta)^2}.
\tag{$\ddagger$}
\]\indent
再注意 (ii) 只在 $V\le R^2$ 时存在边界 $R-\sqrt V\ge0$；而当 $V>R^2$ 时 $c_1=0$ 即已触发 2 补足，(ii) 退化为 $c_1=0$、$\pi_1=\delta V$。此时 1 在“$c_1=0$ 搭便车”与“$c_1=R$ 独自完成”之间比较：
\[
V-R^2\ \ge\ \delta V\iff V\ \ge\ \frac{R^2}{1-\delta}.
\]

令 $V_1=\dfrac{R^2}{(1+\sqrt\delta)^2}$、$V_2=\dfrac{R^2}{1-\delta}$。在各阈值之外，逆向归纳结果为
\[
(c_1^*,c_2^*)=
\begin{cases}
(0,\,0), & V<V_1,\\[5pt]
\bigl(R-\sqrt V,\ \sqrt V\bigr), & V_1<V<R^2,\\[5pt]
(0,\,R), & R^2<V<V_2,\\[5pt]
(R,\,\cdot), & V>V_2.
\end{cases}
\]\indent
阈值处须明确处理无差异：$V=V_1$ 时，$c_1=0,c_2=0$ 与 $c_1=R-\sqrt V,c_2=\sqrt V$ 都可成为均衡路径结果（第二个结果要求合伙人 2 在其无差异时选择补足）；$V=R^2$ 时，合伙人 2 在 $c_1=0$ 无差异，只有选择 $c_2=R$ 才支持合伙人 1 的最优反应；$V=V_2$ 时，合伙人 1 在 $c_1=0$ 与 $c_1=R$ 之间无差异，二者均可为均衡结果。故题干所称“唯一”在这些临界参数上需要补充平局处理规则。
相应收益为
\[
(u_1^*,u_2^*)=
\begin{cases}
(0,\,0), & \text{第 1 种};\\
\bigl(\delta V-(R-\sqrt V)^2,\ 0\bigr), & \text{第 2 种};\\
\bigl(\delta V,\ V-R^2\bigr), & \text{第 3 种};\\
\bigl(V-R^2,\ V\bigr), & \text{第 4 种}.
\end{cases}
\]\indent
当 $\delta\to1$ 时 $\frac{R^2}{1-\delta}\to+\infty$，第 4 种消失；当 $\delta<1$ 时四个区间依次衔接、无空隙。

\myfig{%
\begin{tikzpicture}[scale=0.92, every node/.style={font=\small}]
  \node[pnode] (p1) at (0,0) {};
  \node[above=2pt of p1, nodlab] {$1$};
  \node[pnode] (p2) at (3.4,-3.2) {};
  \node[right=4pt of p2, nodlab] {$2$};
  \node[tnode] (t1) at (-3.4,-3.2) {};
  \node[below=3pt of t1, pay] {$(V-c_1^{2},\ V)$};
  \node[tnode] (t2) at (1.2,-6.0) {};
  \node[below=3pt of t2, pay] {$\bigl(\delta V-c_1^{2},\ V-c_2^{2}\bigr)$};
  \node[tnode] (t3) at (5.6,-6.0) {};
  \node[below=3pt of t3, pay] {$(-c_1^{2},\ -c_2^{2})$};
  \draw[branch] (p1) -- (t1) node[pos=0.5, elab, left=1pt] {$c_1\ge R$};
  \draw[branch] (p1) -- (p2) node[pos=0.55, elab, above left=-2pt] {$c_1<R$};
  \draw[branch] (p2) -- (t2) node[pos=0.55, elab, left=1pt] {$c_1+c_2\ge R$};
  \draw[branch] (p2) -- (t3) node[pos=0.55, elab, right=1pt] {$c_1+c_2<R$};
\end{tikzpicture}%
}{两阶段贡献博弈的博弈树}

\cmt{离开临界值时四种行为依次对应“不启动”“部分贡献并诱使合伙人 2 补足”“由合伙人 2 独自完成”“合伙人 1 独自完成”。临界值处的多重结果或平局处理见上文。}


% ============================================================================
%  2.5  企业专用技术投资与晋升承诺（Prendergast, 1992）
% ============================================================================
\exer{2.5}
假设一个企业希望某工人能投资于企业专门技术 $S$，但此项技术非常模糊，以致法院无法确定工人是否已经掌握。企业无法与职工订立契约偿付工人投资成本；即使职工确实进行了投资，企业也可以声称职工没有进行投资，并且法院无法辨别哪一方是正确的。类似地，如果企业预先支付费用，职工也不能保证会投资于企业的专门技术。但是，企业能够通过（可信的）承诺提职来激励职工进行投资，具体方法如下。假设企业里有两个工作岗位，一个容易（$E$），一个复杂（$D$），并且专门技术对这两个岗位都是有用的，只是对复杂的岗位作用更大一些：$y_{E0}<y_{ES}<y_{DS}<y_{D0}$，其中 $y_{ij}$ 表示职工在岗位 $i$（$i=E$ 或 $D$）工作、技术水平为 $j$（$j=0$ 或 $S$）时的产出，假定企业可以承诺对不同的岗位支付不同的工资 $w_E$ 及 $w_D$，但每一种工资都不低于工人另谋职业的收入，这里我们通过标准化处理，使后者等于 $0$。

博弈进行的时间顺序如下：在时点 0 企业选定 $w_E$ 和 $w_D$，工人观测到企业选择的工资水平。在时点 1 工人加入企业并且能够以成本 $c$ 取得技能 $S$（这里我们忽略了第一阶段的产出及工资，由于工人尚未取得技能，企业聘用工人到岗位 $E$ 是有效率的）。假定 $y_{DS}-y_{E0}>c$ 从而对工人来讲投资学习技术是有效率的。在时点 2 企业观测工人是否取得了技术技能，然后决定是否在工人雇佣的第二（也是最后）阶段提升工人到岗位 $D$。

工人在岗位 $i$ 工作，技术水平为 $j$ 时，企业第二阶段的利润为 $y_{ij}-w_i$。工人第二阶段在岗位 $i$ 工作的收益为 $w_i$ 或 $w_i-c$，决定于工人是否在第一阶段投资于工作技能。请求出博弈的逆向归纳解。参见普伦德加斯特（Prendergast，1992）内容更为丰富的模型。

\sol
时序：时点 0 企业承诺 $(w_E,w_D)$；时点 1 工人选择是否投资 $S$（成本 $c$）；时点 2 企业观测技能状态后决定安置岗位 $E$ 或 $D$。

\textbf{时点 2（企业的最优晋升决策）}。给定工资差 $\Delta w\equiv w_D-w_E$，企业比较两岗位利润：
\[
\text{若工人已掌握技能 }S:\quad \text{晋升到 }D\iff y_{DS}-w_D\ge y_{ES}-w_E
\iff \Delta w\le y_{DS}-y_{ES};
\]\indent
\[
\text{若工人未掌握技能 }0:\quad \text{晋升到 }D\iff y_{D0}-w_D\ge y_{E0}-w_E
\iff \Delta w\le y_{D0}-y_{E0}.
\]\indent
注意由题意 $y_{E0}<y_{ES}<y_{DS}<y_{D0}$，故
\[
y_{D0}-y_{E0} \;>\; y_{DS}-y_{ES}.
\tag{$\dagger$}
\]\indent
即把“未掌握技能者”晋升到 $D$ 的相对吸引力（产出差）反而大于把“已掌握技能者”晋升的相对吸引力。

\textbf{时点 1（工人的投资激励）}。记工人投资后的第二期收益为 $w_D-c$（若被晋升）或 $w_E-c$（若不被晋升）；不投资则为 $w_D$ 或 $w_E$。关键比较是：投资相对不投资额外得到 $-c$（在晋升时）或同样是 $-c$（在不晋升时），即
\[
\text{（投资 payoff）}-\text{（不投资 payoff）}\equiv -c<0
\]\indent
无论企业采用哪种晋升规则都成立——因为晋升决策只取决于岗位而非“是否投资”，工资差 $\Delta w$ 对投资/不投资两类工人一视同仁，投资带来的只是多付出成本 $c$ 而没有对应的工资补偿。因此\textbf{工人永远选择不投资}（$S=0$），与 $(w_E,w_D)$ 无关。

\textbf{时点 0（企业选择工资）}。企业预见到工人不投资，于是只需在“所有工人都只有技能 0”的情形下最大化第二期利润，并满足参与约束 $w_E,w_D\ge 0$。由于 $y_{D0}>y_{E0}$，企业希望把工人安置在岗位 $D$。要让企业在时点 2 选择晋升（即使对未掌握技能者也晋升），只需令 $\Delta w\le y_{D0}-y_{E0}$；同时把工资压到最低 $w_E=w_D=0$ 即可。于是逆向归纳解为：
\[
w_E^*=0,\quad w_D^*=0,\quad \text{工人选择 }S=0,\quad \text{企业将其晋升至岗位 }D.
\]\indent
企业第二期利润为 $y_{D0}$，工人收益为 $0$（等于外部机会成本）。

\cmt{本模型的要点在于 $(\dagger)$ 的不等式：因为“无技能者占据复杂岗位”的产出损失相对更小（与有技能者相比），企业无法用“晋升”把两类工人分离开——只要晋升对无技能者有吸引力，对有技能者必然也有吸引力，而工人投资得不到额外回报。故承诺提职并不能在此参数化下诱致投资；投资之所以“有效率”却无法被实现，正是套牢（hold-up）的体现。}

% ============================================================================
%  2.6  三寡头序贯产量博弈
% ============================================================================
\exer{2.6}
在一个由三个寡头垄断者操纵的市场，反需求函数由 $P(Q)=a-Q$ 给出，这里 $Q=q_1+q_2+q_3$，$q_i$ 表示企业 $i$ 生产的产量。每一企业生产的边际成本为常数 $c$，并且没有固定成本。企业按以下顺序进行产出决策：(1) 企业 1 选择 $q_1\ge 0$；(2) 企业 2 和 3 观测到 $q_1$，并同时分别选择 $q_2$ 和 $q_3$。求出此博弈的子博弈精炼解。

\sol
采用逆向归纳。先求解企业 2、3 在观测到 $q_1$ 后的古诺-纳什均衡，再代回企业 1 的最优化。

\textbf{第 2 阶段（企业 2 与 3 的静态古诺博弈）}。利润函数为
\[
\pi_2=(a-Q-c)q_2=\bigl[a-c-(q_1+q_2+q_3)\bigr]q_2,\qquad
\pi_3=(a-Q-c)q_3=\bigl[a-c-(q_1+q_2+q_3)\bigr]q_3.
\]\indent
一阶条件联立：
\[
\begin{cases}
a-c-q_1-q_3-2q_2=0,\\
a-c-q_1-q_2-2q_3=0,
\end{cases}
\quad\Longrightarrow\quad
q_2^*(q_1)=q_3^*(q_1)=\dfrac{a-c-q_1}{3}.
\]

\textbf{第 1 阶段（企业 1 的最优化）}。将上式代入，企业 1 面对
\[
\pi_1=(a-Q-c)q_1
=\Bigl(a-c-q_1-2\cdot\frac{a-c-q_1}{3}\Bigr)q_1
=\frac{1}{3}(a-c-q_1)q_1.
\]\indent
一阶条件：
\[
\frac{d\pi_1}{dq_1}=\frac{1}{3}(a-c-2q_1)=0
\quad\Longrightarrow\quad
q_1^*=\frac{a-c}{2}.
\]\indent
从而
\[
q_2^*=q_3^*=\frac{a-c}{6}.
\]

\textbf{子博弈精炼纳什均衡（SPNE）}：
\[
\boxed{\,q_1^*=\dfrac{a-c}{2},\quad q_2^*=q_3^*=\dfrac{a-c}{6}\,}.
\]\indent
总产出 $Q^*=\frac{5}{6}(a-c)$，价格 $P^*=a-Q^*=\frac{1}{6}(a+5c)$。\cmt{先动者（企业 1）凭借“先行优势”产量最高、利润最高，这正是序贯古诺（Stackelberg 型）三寡头结构的结论。}

% ============================================================================
%  2.7  垄断工会与 $n$ 个企业
% ============================================================================
\exer{2.7}
假设一个工会是一个寡头垄断市场中所有企业唯一的劳动力供给者，就像汽车工人联合会（UAW）对于通用、福特、克莱斯勒等大型汽车厂家。令博弈各方行动的时间顺序类似于第 2.1.C 节的模型：(1) 工会确定单一的工资要求 $w$，适用于所有的企业；(2) 每家企业 $i$ 了解到（并接受）$w$，然后同时分别选择各自的雇佣水平 $L_i$；(3) 工会的收益为 $(w-w_a)L$，其中 $w_a$ 为工会成员到另外行业谋职可取得的收入，$L=L_1+\cdots+L_n$ 为工会在本行业企业的总就业水平；企业 $i$ 的利润为 $\pi_i(w,L_i)$，其中决定企业 $i$ 利润水平的要素如下。所有企业都有同样的生产函数：产出等于劳动力；$q_i=L_i$。市场总产出为 $Q=q_1+\cdots+q_n$ 时的市场出清价格为 $p(Q)=a-Q$。为使问题简化，假设企业除工资支出外没有另外的成本。请求出此博弈的子博弈精炼解。在子博弈精炼解中，企业的数量是如何（以及为什么）影响工会效用的？

\sol
时序：第 1 阶段工会定 $w$；第 2 阶段各企业在给定 $w$ 下同时选择 $L_i$。

\textbf{第 2 阶段（给定 $w$ 下企业的古诺博弈）}。企业 $i$ 利润
\[
\pi_i=(p(Q)-w)L_i=(a-Q-w)L_i=(a-w-L)L_i,
\]\indent
其中 $L=\sum_j L_j$。一阶条件：
\[
\frac{\partial\pi_i}{\partial L_i}=a-w-L-L_i=0
\quad\Longrightarrow\quad
L_i=\frac{a-w-L}{1}\,\Rightarrow\,\text{对称均衡 }L_i^*=\frac{L^*}{n}.
\]\indent
代入 $a-w-L^*-L^*/n=0$ 得
\[
L^*=L(w)=\frac{n(a-w)}{n+1}.
\]

\textbf{第 1 阶段（工会的最优化）}。工会最大化
\[
U(w)=(w-w_a)L(w)=(w-w_a)\frac{n(a-w)}{n+1}.
\]\indent
一阶条件：
\[
\frac{dU}{dw}=\frac{n}{n+1}\bigl[(a-w)-(w-w_a)\bigr]=0
\quad\Longrightarrow\quad
2w=a+w_a\;\Longrightarrow\; w^*=\frac{a+w_a}{2}.
\]\indent
于是
\[
L^*=\frac{n(a-w^*)}{n+1}
=\frac{n\bigl(a-\frac{a+w_a}{2}\bigr)}{n+1}
=\frac{n(a-w_a)}{2(n+1)},
\]\indent
工会效用
\[
U^*=(w^*-w_a)L^*
=\frac{a-w_a}{2}\cdot\frac{n(a-w_a)}{2(n+1)}
=\frac{n(a-w_a)^2}{4(n+1)}.
\]

\textbf{企业数量的影响}。注意 $n/(n+1)$ 随 $n$ 严格递增，故
\[
\frac{\partial U^*}{\partial n}>0.
\]\indent
即\textbf{行业内企业越多，工会的效用越高}。原因：工资被钉在 $w^*=(a+w_a)/2$ 上与企业数量无关，但总就业 $L^*$ 随 $n$ 而增加（因为更多企业竞争劳动力使就业对工资更富弹性），从而工会以固定工资“卖出”了更多劳动，效用上升。\cmt{这与直觉相反——通常认为买方垄断力量随雇主减少而增强，但此处工会是垄断卖方且工资由“租金抽取”决定，企业数目增加使劳动力需求更富弹性，反而有利于工会。}

% ============================================================================
%  2.8  引入破坏的工作竞赛（Lazear, 1989）
% ============================================================================
\exer{2.8}
根据拉齐尔（1989），对第 2.2.D 节的工作竞赛模型修改如下：令工人 $i$ 的产出为 $y_i=e_i-\frac12 s_i+\varepsilon_i$，其中 $s_i\ge 0$ 代表了工人的故意破坏；并且工人 $i$ 由于（生产及破坏）付出努力产生的负效用为 $g(e_i)+g(s_i)$。证明最优的奖励水平 $w_H-w_L$ 要小于（书中）不考虑进行破坏的情况。

\sol
竞赛结构（沿用第 2.2.D 节）：两工人同时选择行动；相对产出高者得 $w_H$，低者得 $w_L$（$w_H>w_L$）。记噪声差 $\eta=\varepsilon_i-\varepsilon_j$ 的密度为 $f$，在对称处对胜率关于“相对产出”的边际效应记为 $\beta>0$（即工人 $i$ 相对产出每提高一单位，其胜率约提高 $\beta$；由对称性 $\beta=f(0)/2>0$）。

\textbf{胜率对行动的导数}。工人 $i$ 的相对产出为
\[
y_i-y_j=(e_i-e_j)-\tfrac12(s_i-s_j)+(\varepsilon_i-\varepsilon_j).
\]\indent
故
\[
\frac{\partial\Pr(\text{赢})}{\partial e_i}=+\beta,\qquad
\frac{\partial\Pr(\text{赢})}{\partial s_i}=+\frac{\beta}{2}.
\]\indent
（破坏 $s_i$ 出现在工人 $i$ 自身产出式中以 $-\tfrac12 s_i$ 出现，等价于降低自身相对产出，因此这里采用教材所给形式时 $s_i$ 降低自身产出；但为使“破坏”能发生，其动机必须来自伤害对手——以下给出更具标准性的等价表述：把 $s_i$ 理解为“对对手的破坏”，它提高 $i$ 的相对产出，故 $\partial\Pr/\partial s_i=+\beta$，且破坏成本同为 $g(s_i)$。两种表述的定性结论一致，下面采用后者，因其能导出“破坏确实发生”这一事实。）

\textbf{工人的最优化（对称均衡）}。给定工资差 $\Delta w\equiv w_H-w_L$，工人 $i$ 选择 $e_i,s_i$ 最大化
\[
\max_{e_i,s_i\ge 0}\ \Delta w\cdot\Pr_i(e_i,s_i)-g(e_i)-g(s_i).
\]\indent
一阶条件（内点）：
\[
g'(e_i)=\Delta w\cdot\beta,\qquad g'(s_i)=\Delta w\cdot\beta.
\]\indent
因 $g(\cdot)$ 相同且严格凸，得 $e_i^*=s_i^*\equiv x(\Delta w)$，且
\[
g'\bigl(x(\Delta w)\bigr)=\Delta w\cdot\beta.\tag{1}
\]\indent
可见提高 $\Delta w$ 会同时等比例地提高努力与破坏。

\textbf{企业的目标}。企业每单位劳动净产出（相对自身努力减去所受破坏）的期望为 $e_i-\mathbb{E}[s_j]=x-x=0$（对称下努力与所受破坏抵消），而企业却要支付奖金 $w_H$（赢家）与 $w_L$（输家）。因此\textbf{提高 $\Delta w$ 对净产出的边际贡献被“努力”与“破坏”的相互抵消而削弱}：在不存在破坏的书中模型里，$\partial(\text{净产出})/\partial(\Delta w)>0$ 完全由努力 $e$ 贡献；加入破坏后，增加的努力被等量的破坏抵消，净产出对 $\Delta w$ 的敏感度下降。

更正式地，企业选择 $\Delta w$ 最大化（每工人）期望利润
\[
\Pi= \mathbb{E}[y_i]-(\text{期望工资})
= \bigl(e^*-\mathbb{E}[s_j]\bigr)-\bigl[\tfrac12 w_H+\tfrac12 w_L\bigr].
\]\indent
无破坏时 $s\equiv 0$，$e^*$ 随 $\Delta w$ 上升使 $\Pi$ 在某一正 $\Delta w^*$ 处最大化；有破坏时 $e^*=\mathbb{E}[s_j]$，净产出项恒为 $0$，企业仅剩压低工资的压力，故最优 $\Delta w$ 必须低于无破坏情形。由 $(1)$ 知 $\Delta w$ 上升只会推高 $x$ 从而同时推高努力与破坏的浪费性成本，而不会带来净产出收益，于是
\[
\boxed{\Delta w_{\text{有破坏}}^{*}<\Delta w_{\text{无破坏}}^{*}}.
\]\indent
\cmt{结论的经济含义：破坏是一种“零和且浪费”的活动，较大的奖金差距在激励努力的同时也激励破坏，而破坏不创造任何产出，故企业会刻意把奖金差距定得比无破坏时更小，以抑制破坏。}

% ============================================================================
%  2.9  关税取消前后的工资—就业博弈
% ============================================================================
\exer{2.9}
考虑两个国家，在时点 1，两国的关税都非常高，以致根本没有贸易。在每一个国家内，工资和就业的决定与第 2.1.C 节垄断工会模型相同。在时点 2，所有的关税都取消了，每一工会决定本国的工资水平，但每一企业同时为两国市场生产产品。假定每一国家的反需求函数为 $P(Q)=a-Q$，其中 $Q$ 为该国市场上的总产量。令每一企业的生产函数为 $q=L$，从而工资为企业的惟一成本，并且令工会的效用函数为 $U(w,L)=(w-w_0)L$，其中 $w_0$ 为工人另谋职业可得收入。请解出时点 1 博弈的逆向归纳解。下面考虑时点 2 的博弈。第一，两工会同时选择工资水平 $w_1$ 和 $w_2$，然后企业 $i$ 观测到工资并选择其国内、国外市场的产出水平，企业 $i$ 的产出（内销和出口）分别用 $h_i$ 和 $e_i$ 表示。企业 $i$ 所有的生产都在本国，于是总成本为 $w_i(h_i+e_i)$，解出子博弈精炼解。证明工资、就业和利润（并且由此决定的工会效用和消费者剩余）在关税取消后都提高了。参见 Huizinga（1989）同类问题的其他例子。

\sol
两个对称国家，记两国为 $i,j\in\{1,2\}$。

\textbf{时点 1：高关税，无贸易}。两国相互独立且同质。企业 $i$ 利润
\[
\pi_i=(a-L_i)L_i-w_iL_i,
\]\indent
一阶条件给出劳动需求 $L_i(w_i)=\dfrac{a-w_i}{2}$。工会效用
\[
U_i=(w_i-w_0)L_i(w_i)=\frac{(w_i-w_0)(a-w_i)}{2}.
\]\indent
一阶条件得 $w_i^*=\dfrac{a+w_0}{2}$，从而
\[
L_i^*=\frac{a-w_0}{4},\qquad
\pi_i^*=\frac{(a-w_0)^2}{16}.
\]\indent
逆向归纳解：
\[
\Bigl\{w_i^*=w_j^*=\frac{a+w_0}{2},\ L_i^*=L_j^*=\frac{a-w_0}{4},\ \pi_i^*=\pi_j^*=\frac{(a-w_0)^2}{16}\Bigr\}.
\]

\textbf{时点 2：关税取消，两国贸易}。企业 $i$ 在两国市场销售，国内销量 $h_i$、出口 $e_i$；企业 $j$ 对应为 $h_j,e_j$。企业总成本 $w_i(h_i+e_i)$。利润：
\[
\pi_i=h_i(a-h_i-e_j)+e_i(a-h_j-e_i)-w_i(h_i+e_i).
\]\indent
对 $h_i,e_i$ 分别求一阶条件：
\[
\frac{\partial\pi_i}{\partial h_i}=a-2h_i-e_j-w_i=0,\qquad
\frac{\partial\pi_i}{\partial e_i}=a-2e_i-h_j-w_i=0.
\]\indent
由对称性解得（联立两国四个方程）：
\[
h_i^*=e_i^*=\frac{a+w_j-2w_i}{3}.
\]\indent
故企业 $i$ 的劳动力总需求
\[
L_i^*=h_i^*+e_i^*=\frac{2(a+w_j-2w_i)}{3}.
\]\indent
工会效用
\[
U_i=(w_i-w_0)L_i^*
=\frac{2(w_i-w_0)(a+w_j-2w_i)}{3}.
\]\indent
对称均衡下 $w_1^*=w_2^*=w^*$，一阶条件：
\[
\frac{\partial U_i}{\partial w_i}=\frac{2}{3}\bigl[(a+w_j-2w_i)-2(w_i-w_0)\bigr]=0
\;\Longrightarrow\;
a+w_j+2w_0-4w_i=0.
\]\indent
代入 $w_j=w_i$ 得
\[
w^*=\frac{a+2w_0}{3}.
\]\indent
于是
\[
L^*=\frac{4(a-w_0)}{9},\qquad
\pi_i^*=h_i^*(a-h_i^*-e_j^*)+e_i^*(a-h_j^*-e_i^*)-w_i^*L_i^*
=\frac{8(a-w_0)^2}{81}.
\]

\textbf{比较时点 1 与时点 2}：
\[
w:\ \frac{a+2w_0}{3}>\frac{a+w_0}{2}\quad(\text{因 }a>w_0);
\]\indent
\[
L:\ \frac{4(a-w_0)}{9}>\frac{a-w_0}{4};
\]\indent
\[
\pi:\ \frac{8(a-w_0)^2}{81}>\frac{(a-w_0)^2}{16}\quad\bigl(8/81\approx0.099>1/16=0.0625\bigr).
\]\indent
工资、就业、企业利润均上升；由利润上升与价格下降（贸易扩大供给）可知消费者剩余亦上升，工会效用随之上升。$\blacksquare$


% ============================================================================
%  2.10  两阶段重复博弈中的 SPNE（x 的阈值）
% ============================================================================
\exer{2.10}
下图所示的同时行动博弈重复进行两次，并且第二阶段开始前双方均可观测到第一阶段的结果，不考虑贴现因素。变量 $x$ 大于 $4$，因而 $(4,4)$ 在一次性博弈中并不是均衡收益。对什么样的 $x$，（双方参与者同时采取）下述战略是一个子博弈精炼纳什均衡？

第一阶段选择 $Q$，如果第一阶段的结果为 $(Q_1,Q_2)$，在第二阶段选择 $P$；如果第一阶段的结果为 $(y,Q_2)$，其中 $y\neq Q_1$，第二阶段选择 $R$；如果第一阶段的结果为 $(Q_1,z)$，其中 $z\neq Q_2$，第二阶段选择 $S$；如果第一阶段结果为 $(y,z)$，其中 $y\neq Q_1$ 且 $z\neq Q_2$，则在第二阶段选择 $P$。

\sol
一次性阶段博弈的收益矩阵（行=参与者 1 的行动，列=参与者 2 的行动）为
\begin{paytab}
\begin{tabular}{ccccc}
 & \pthead{$P_2$} & \pthead{$Q_2$} & \pthead{$R_2$} & \pthead{$S_2$}\\[2pt]
\mapline
$P_1$ & $2,2$ & $x,0$ & $-1,0$ & $0,0$\\[4pt]
\mapline
$Q_1$ & $0,x$ & $4,4$ & $-1,0$ & $0,0$\\[4pt]
\mapline
$R_1$ & $0,0$ & $0,0$ & $0,2$ & $0,0$\\[4pt]
\mapline
$S_1$ & $0,-1$ & $0,-1$ & $-1,-1$ & $2,0$\\[4pt]
\mapline
\end{tabular}
\end{paytab}

题中战略（双方对称采用）可写为：参与者 $i$ 在第 1 阶段选 $Q_i$；在第 2 阶段，若历史为 $(Q_1,Q_2)$ 或 $(y,z)$（双方都偏离）则选 $P_i$；若自己偏离、对手选 $Q$ 即历史为 $(y,Q_2)$ 则选 $R_i$；若自己选 $Q$、对手偏离即历史为 $(Q_1,z)$ 则选 $S_i$。记该策略组合为 $\sigma$。

\textbf{（一）检验各第 2 阶段子博弈中 $\sigma$ 是否构成阶段博弈 NE}。
\begin{enumerate}[leftmargin=*]
  \item 历史 $(Q_1,Q_2)$：第 2 阶段行动 $(P_1,P_2)$，收益 $(2,2)$。对参与者 1，在 $P_2$ 下各行动收益为 $2,0,0,0$，最优为 $P_1$；对参与者 2，在 $P_1$ 下各列收益为 $2,0,0,0$，最优为 $P_2$。故 $(P,P)$ 是 NE。
  \item 历史 $(y,Q_2)$（如 $(P_1,Q_2)$）：第 2 阶段行动 $(R_1,R_2)$，收益 $(0,2)$。在 $R_2$ 下列收益为 $-1,-1,0,0$，最优为 $R_1$；在 $R_1$ 下行收益为 $0,0,2,0$，最优为 $R_2$。故 $(R,R)$ 是 NE。
  \item 历史 $(Q_1,z)$（如 $(Q_1,P_2)$）：第 2 阶段行动 $(S_1,S_2)$，收益 $(2,0)$。在 $S_2$ 下列收益为 $0,0,0,2$，最优为 $S_1$；在 $S_1$ 下行收益为 $-1,-1,-1,0$，最优为 $S_2$。故 $(S,S)$ 是 NE。
  \item 历史 $(y,z)$（双方偏离）：第 2 阶段 $(P_1,P_2)=(2,2)$，同前是 NE。
\end{enumerate}
因此 $\sigma$ 在\textit{所有}第 2 阶段子博弈中都给出阶段博弈的 NE。

\textbf{（二）检验第 1 阶段无有利可图的单方面偏离}。若不偏离，双方第 1 阶段 $(Q_1,Q_2)=(4,4)$，第 2 阶段 $(P,P)=(2,2)$，总收益均为 $6$。

考虑参与者 1 偏离到 $y\neq Q_1$（参与者 2 仍按 $\sigma$ 选 $Q_2$）：
\begin{enumerate}[leftmargin=*]
  \item 若选 $P_1$：第 1 阶段 $(P_1,Q_2)=(x,0)$，参与者 1 得 $x$；历史 $(P_1,Q_2)$ 触发第 2 阶段 $(R_1,R_2)$，参与者 1 得 $0$。总收益 $x$。
  \item 若选 $R_1$：第 1 阶段 $(-1,0)$，第 2 阶段 $(R,R)$ 得 $0$，$(R_1,Q_2)=(0,0)$，第 1 阶段得 $0$；历史 $(R_1,Q_2)$ 属于 $y\neq Q_1$，触发第 2 阶段 $(R_1,R_2)=(0,2)$，参与者 1 得 $0$。总 $0$。
  \item 若选 $S_1$：第 1 阶段 $(0,-1)$，第 2 阶段得 $0$，总 $0$。
\end{enumerate}
参与者 1 偏离到 $P_1$ 可获 $x$，当且仅当 $x>6$ 时有利可图。

同理，参与者 2 偏离到 $z\neq Q_2$（参与者 1 仍选 $Q_1$）：
\begin{enumerate}[leftmargin=*]
  \item 选 $P_2$：历史 $(Q_1,P_2)$，第 1 阶段参与者 2 得 $x$，第 2 阶段 $(S,S)$ 得 $0$，总 $x$。
  \item 选 $R_2$ 或 $S_2$：第 1 阶段均得 $0$（见 $(Q_1,R_2)=(-1,0)$、$(Q_1,S_2)=(0,0)$），第 2 阶段得 $0$，总 $0$。
\end{enumerate}
参与者 2 当且仅当 $x>6$ 时有利可图。

\textbf{结论}。所给战略是子博弈精炼纳什均衡，当且仅当
\[
\boxed{4<x\le 6}.
\]\indent
（当 $x=6$ 时偏离与不偏离总收益均为 $6$，属弱无差异，不构成本质性偏离；故阈值取 $x\le 6$。）

% ============================================================================
%  2.11  两阶段重复博弈中能否在第一阶段达到 (4,4)
% ============================================================================
\exer{2.11}
下图的同时行动博弈重复进行两次，第二阶段开始前参与双方可以观测到第一阶段的结果。不考虑贴现因素。在一个纯战略子博弈精炼纳什均衡中，能否在第一阶段达到 $(4,4)$ 的收益？如果可以，给出相应的战略。如果不能，证明为什么。

\sol
一次性阶段博弈的收益矩阵为
\begin{paytab}
\begin{tabular}{ccccc}
 & \pthead{$L$} & \pthead{$C$} & \pthead{$R$}\\[2pt]
\mapline
$T$ & $3,1$ & $0,0$ & $5,0$\\[4pt]
\mapline
$M$ & $2,1$ & $1,2$ & $3,1$\\[4pt]
\mapline
$B$ & $1,2$ & $0,1$ & $4,4$\\[4pt]
\mapline
\end{tabular}
\end{paytab}

\textbf{阶段博弈的纯战略 NE}：用划线法可知纯战略 NE 只有两个：$(T,L)$ 与 $(M,C)$。注意 $(B,R)$ \textbf{并不是}阶段博弈的 NE：给定参与者 2 选 $R$，参与者 1 选 $T$ 可得 $5>4$，有偏离激励（尽管 $(4,4)$ 是全矩阵的最高对称收益，却不是均衡收益）。

\textbf{构造一个纯战略 SPNE，使第 1 阶段达到 $(4,4)$}。令双方采用如下触发战略 $\sigma$：
\begin{quote}
第 1 阶段选 $(B,R)$；若第 1 阶段结果恰为 $(B,R)$，第 2 阶段选 $(T,L)$；若第 1 阶段为任何其他结果，第 2 阶段选 $(M,C)$。
\end{quote}
检验：
\begin{enumerate}[leftmargin=*,label=(\roman*)]
  \item 第 2 阶段：无论第 1 阶段出现何种历史，战略规定的行动组合或为 $(T,L)$、或为 $(M,C)$，而二者均为阶段博弈的纯 NE，故每个第 2 阶段子博弈中都构成 NE，满足子博弈精炼要求。
  \item 第 1 阶段：若双方遵循 $\sigma$，第 1 阶段 $(B,R)=(4,4)$，第 2 阶段 $(T,L)=(3,1)$，总收益为 $(7,5)$。

  参与者 1 单方面偏离（偏离触发第 2 阶段 $(M,C)$，参与者 2 阶段得 $1$）：偏离到 $T$ 得 $5+1=6$，到 $M$ 得 $3+1=4$，均小于 $7$；故不偏离。

  参与者 2 单方面偏离（偏离触发第 2 阶段 $(M,C)$，参与者 2 阶段得 $2$）：偏离到 $L$ 得 $2+2=4$，到 $C$ 得 $1+2=3$，均小于 $5$；故不偏离。
\end{enumerate}
因此 $\sigma$ 是纯战略 SPNE，且第 1 阶段达到 $(4,4)$。

\cmt{结论：可以。由于 $(B,R)$ 本身不是阶段博弈的 NE，维持它必须依靠第 2 阶段的触发惩罚：一旦第 1 阶段偏离，第 2 阶段便转投另一 NE $(M,C)$，使偏离者（尤其参与者 1）损失 $3-1=2$ 的第 2 阶段收益，足以抵消其当期偏离所得 $+1$。任何第 2 阶段被规定为某个阶段博弈 NE 的触发式战略，只要满足上述两步偏离检验，都可构成纯战略 SPNE。}

% ============================================================================
%  2.12  概念：战略 / 子博弈 / SPNE
% ============================================================================
\exer{2.12}
何为重复博弈的战略？何为重复博弈的子博弈？何为子博弈精炼纳什均衡？

\sol
\begin{enumerate}[leftmargin=*,label=(\roman*)]
  \item \textbf{重复博弈的战略}。在基于阶段博弈 $G$ 的（有限或无限）重复博弈中，一个参与者的战略是一份“完整计划”，它针对\textit{每一条可能出现的历史（history）}规定该参与者下一阶段应采取的阶段博弈行动。历史指到该时刻为止所有参与者过去各阶段行动的有序记录；对第 1 阶段，历史为空。因此重复博弈的战略远多于阶段博弈的行动，因为它必须说明“在任何可能的过去轨迹下”怎么做。
  \item \textbf{重复博弈的子博弈}。从一个“信息完备的历史节点”开始的后续博弈构成一个子博弈；对重复博弈而言，由于每一阶段结束都观测到全部行动，\textit{每一个历史之后}都对应一个子博弈（对于有限重复，连同初始节点共 $1+$ 历史数 个子博弈；无限重复则有可数无穷多个）。每个这样的子博弈本身也是一个（从当前历史起）重复博弈。
  \item \textbf{子博弈精炼纳什均衡（SPNE）}。一个战略组合如果：\textbf{（a）}它是整个博弈的纳什均衡；且 \textbf{（b）}它在\textit{每一个}子博弈（即针对每一条历史）中都给出该子博弈的纳什均衡，则称为子博弈精炼纳什均衡。换言之，SPNE 要求没有任何参与者在任何信息集（包括非初始的）上有偏离激励——它剔除了那些依靠“不可置信的威胁”来维持的纳什均衡。
\end{enumerate}


% ============================================================================
%  2.13  贝特兰德双头垄断的垄断价格维持（grim trigger）
% ============================================================================
\exer{2.13}
回顾习题 1.7 中静态贝特兰德双头垄断模型（同类产品）：企业同时给产品定价；对企业 $i$ 产品的需求，如果 $p_i<p_j$，为 $a-p_i$；如果 $p_i>p_j$，为 $0$；如果 $p_i=p_j$，则为 $(a-p_i)/2$；边际成本为 $c<a$。考虑基于上述阶段博弈的无限重复博弈。证明在一个子博弈精炼纳什均衡中，当且仅当 $\delta>1/2$ 时，企业可以使用触发战略（只要发生任何背离，就永远转向阶段博弈的纳什均衡）以维持垄断条件下的价格水平。

\sol
\textbf{阶段博弈回顾}。同质产品贝特兰德竞争的唯一 NE 是 $p_1=p_2=c$，利润均为 $0$。若双方都定垄断价格 $p^m$，则须满足
\[
\max_p (p-c)(a-p)\quad\Longrightarrow\quad p^m=\frac{a+c}{2},
\]\indent
此时每家企业利润（平分市场）为
\[
\pi^m=\frac{(p^m-c)(a-p^m)}{2}=\frac{(a-c)^2}{8}.
\]

\textbf{触发战略}。双方约定：始终定价 $p^m$；一旦任一方偏离，则永远转为阶段博弈 NE（价格战，$p=c$，双方利润 $0$）。

\textbf{单期偏离收益}。若企业 1 在某期把价格略微降到 $p^m-\varepsilon$ 以下以独享整个市场，其当期最大利润为
\[
\pi^D=\max_p (p-c)(a-p)=\frac{(a-c)^2}{4},
\]\indent
之后各期因触发战略而永久得到 $0$。

\textbf{不偏离的现值}。若始终合作，每期利润为 $\pi^m$，无限期折现和为
\[
V=\frac{\pi^m}{1-\delta}=\frac{(a-c)^2}{8(1-\delta)}.
\]

\textbf{激励相容（IC）条件}：
\[
\begin{aligned}
V\ge \pi^D
&\iff \frac{(a-c)^2}{8(1-\delta)}\ge \frac{(a-c)^2}{4}\\
&\iff \frac{1}{8(1-\delta)}\ge \frac14\\
&\iff \frac{1}{1-\delta}\ge 2\\
&\iff 1-\delta\le \frac12\\
&\iff \delta\ge \frac12.
\end{aligned}
\]\indent
由于“只要发生任何背离就永远转向 NE”的威胁在 $\delta>1/2$ 时可置信（偏离者未来损失大于当期收益），故

\[
\boxed{\text{当且仅当 }\delta>1/2\text{ 时，可用触发战略维持垄断价格。}}
\]

% ============================================================================
%  2.14  需求波动下的贝特兰德 grim trigger
% ============================================================================
\exer{2.14}
假设在习题 2.13 中的无限重复贝特兰德博弈中的需求随机波动：在每一阶段，需求函数的截距是 $a_H$ 的概率为 $\pi$；是 $a_L$（$a_L<a_H$）的概率为 $1-\pi$；不同阶段间的需求相互独立。假设每一阶段两企业在选择当期价格之前都可以观测到本期的需求水平。两种需求水平下的垄断价格 $p_H$ 和 $p_L$ 分别是多少？解出 $\delta^*$，使得在子博弈精炼纳什均衡中，两企业能够采用触发战略维持上述垄断价格水平（即对 $i=H,L$，需求为 $a_i$ 时选择价格 $p_i$）的最小贴现因子 $\delta$ 的值。并对处于 $1/2$ 和 $\delta^*$ 之间的所有 $\delta$，求出最高价格 $p(\delta)$，使得在子博弈精炼纳什均衡中，企业可通过运用触发战略，在高需求时维持价格 $p(\delta)$，在低需求时维持价格 $p_L$。

\sol
\textbf{垄断价格}。对截距为 $a_i$ 的需求，垄断价格（平分市场）为
\[
p_i^m=\frac{a_i+c}{2},\qquad i=H,L,
\]\indent
对应每期单方利润
\[
\pi_i^m=\frac{(a_i-c)^2}{8},\quad i=H,L.
\]

\textbf{触发战略与 $\delta^*$}。约定：每期观测到需求后，高需求定 $p_H^m$、低需求定 $p_L^m$；任一方偏离则永久转为价格战（$p=c$，利润 $0$）。记
\[
H\equiv \frac{(a_H-c)^2}{8},\qquad L\equiv \frac{(a_L-c)^2}{8},\qquad
S\equiv \pi H+(1-\pi)L.
\]\indent
合作下的期望现值（每期初）为
\[
V=\frac{S}{1-\delta}.
\]\indent
在高需求状态，合作当期得 $H$，偏离则当期的全市场利润为 $2H=\frac{(a_H-c)^2}{4}$（略低于 $p_H^m$ 即可独享），之后永久为 $0$。激励相容要求
\[
H+\delta V\ge 2H\quad\Longleftrightarrow\quad \delta V\ge H.
\]\indent
代入 $V=S/(1-\delta)$：
\[
\frac{\delta S}{1-\delta}\ge H
\;\Longrightarrow\;
\delta\ge \frac{H}{S+H}
=\frac{(a_H-c)^2}{(1+\pi)(a_H-c)^2+(1-\pi)(a_L-c)^2}.
\]\indent
低需求状态的 IC 约束为 $\delta V\ge L$，因 $H>L$ 自动满足。故最小贴现因子
\[
\boxed{\delta^*=\frac{(a_H-c)^2}{(1+\pi)(a_H-c)^2+(1-\pi)(a_L-c)^2}}.
\]\indent
（可见 $\delta^*<\frac{H}{H+(1-\pi)L}<\frac12$？注意此处因 $S$ 含期望，通常 $\delta^*<1/2$ 亦可；但仅当高需求时单独维持 $p_H^m$ 才需 $\delta>1/2$，见下。）

\textbf{对 $\delta\in(1/2,\delta^*)$ 的最高高需求价格 $p(\delta)$}。此时无法维持完整的 $p_H^m$，但可在高需求维持某个较低价格 $p\le p_H^m$，同时在低需求仍维持 $p_L^m$。令
\[
H(p)\equiv\frac{(p-c)(a_H-p)}{2},\qquad
L_0\equiv\frac{(p_L^m-c)(a_L-p_L^m)}{2}=L.
\]\indent
合作的期望现值 $V(p)=[\pi H(p)+(1-\pi)L_0]/(1-\delta)$。高需求下偏离可得全市场利润 $(p-c)(a_H-p)=2H(p)$。IC 条件：
\[
H(p)+\delta V(p)\ge 2H(p)\quad\Longleftrightarrow\quad \delta V(p)\ge H(p).
\]\indent
即
\[
\frac{\delta\bigl[\pi H(p)+(1-\pi)L_0\bigr]}{1-\delta}\ge H(p)
\;\Longrightarrow\;
H(p)\ge \frac{\delta(1-\pi)L_0}{1-\delta(1+\pi)}.
\]\indent
（需 $\delta<1/(1+\pi)$，对题设区间成立。）因 $H(p)$ 在 $p\le p_H^m$ 上递增，最高可持续价格 $p(\delta)$ 由等号决定：
\[
\frac{(p-c)(a_H-p)}{2}=\frac{\delta(1-\pi)L_0}{1-\delta(1+\pi)}.
\]\indent
解之（取不超过 $p_H^m$ 的根）：
\[
\boxed{
p(\delta)=\frac{a_H+c-\sqrt{(a_H-c)^2-8\,\dfrac{\delta(1-\pi)L_0}{1-\delta(1+\pi)}}}{2}
}.
\]\indent
当 $\delta\uparrow\delta^*$ 时右端根号趋于 $(a_H-c)^2-8L$（即回到 $H(p_H^m)=H$ 的情形），得 $p(\delta^*)=p_H^m$；当 $\delta\downarrow 1/2$ 时 $p(\delta)$ 下降。\cmt{直观：需求波动使平均合作收益 $S$ 低于纯高需求情形，故维持最高价所需的 $\delta^*$ 实际上低于 $1/2$ 即可；但对“高需求单独维持 $p_H^m$”则需要 $\delta>1/2$，题设区间 $(1/2,\delta^*)$ 给出的是在还能维持低需求垄断价前提下、高需求价可上浮到的范围。}

% ============================================================================
%  2.15  $n$ 个企业的古诺垄断产出维持
% ============================================================================
\exer{2.15}
假设在一个古诺寡头垄断中有 $n$ 个企业。反需求函数由 $P(Q)=a-Q$ 给出，其中 $Q=q_1+\cdots+q_n$。考虑基于这一阶段博弈的无限重复博弈。求出最小的 $\delta$ 值，使得在子博弈精炼纳什均衡中企业能够运用触发战略维持垄断产出水平。这一结果随 $n$ 的变化将如何变化？为什么？如果 $\delta$ 非常之小，以致企业无法通过采用触发战略来维持垄断产出，求出能够通过触发战略维持住的、使利润最大化的相应的子博弈精炼纳什均衡。

\sol
设各企业对称地维持每期产量 $q$。垄断总产出 $Q^m=a/2$，故每家企业垄断产量 $q^m=a/(2n)$，单方利润
\[
\pi^m=\bigl(a-nq^m\bigr)q^m=\frac{a^2}{4n}.
\]

\textbf{偏离收益}。给定其余 $n-1$ 家均产 $q^m$，某家偏离选择 $q^D$ 最大化 $(a-q^D-(n-1)q^m)q^D$，一阶条件给出
\[
q^D=\frac{a+(n-1)q^m}{2}\cdot? \quad\text{由 }a-2q^D-(n-1)q^m=0
\;\Longrightarrow\;
q^D=\frac{a-(n-1)q^m}{2}=\frac{a(n+1)}{4n}.
\]\indent
偏离当期利润
\[
\pi^D=(a-q^D-(n-1)q^m)q^D=(q^D)^2=\frac{a^2(n+1)^2}{16n^2}.
\]

\textbf{惩罚}。偏离后触发战略转入阶段博弈古诺 NE：每家企业 $q^c=a/(n+1)$，利润 $\pi^c=(a/(n+1))^2$。

\textbf{IC 条件（维持垄断）}：
\[
\frac{\pi^m}{1-\delta}\ge \pi^D+\frac{\delta}{1-\delta}\pi^c.
\]\indent
代入并消去 $a^2$：
\[
\frac{1}{4n(1-\delta)}\ge \frac{(n+1)^2}{16n^2}+\frac{\delta}{(1-\delta)(n+1)^2}.
\]\indent
整理得
\[
\delta\ge \frac{(n+1)^2}{(n+1)^2+4n}\equiv \delta_{\min}(n).
\]\indent
因此最小贴现因子为
\[
\boxed{\delta_{\min}(n)=\frac{(n+1)^2}{(n+1)^2+4n}}.
\]

\textbf{随 $n$ 的变化}。$\delta_{\min}(n)$ 随 $n$ 严格递增，且 $\delta_{\min}(n)\to 1$（当 $n\to\infty$）。原因：企业越多，单方面偏离所攫取的“超额市场份额”越小（因 $q^D$ 相对总产出占比下降），但同时惩罚（回到古诺竞争）对每家企业的损失也随 $n$ 增大而相对变小——净效果是维持合谋所需的可置信惩罚更难（因为竞争 NE 本身利润随 $n$ 减小得不快），故所需 $\delta$ 升高。特别地 $n=2$ 时 $\delta_{\min}=9/17$。

\textbf{$\delta$ 很小时的最优可持续对称均衡}。令各企业对称地维持产量 $q$（总产出 $Q=nq$），合作利润 $\pi(q)=(a-nq)q$，偏离利润 $\pi^D(q)=\bigl[\frac{a-(n-1)q}{2}\bigr]^2$（因偏离者占有的最优产量 $q^D=\frac{a-(n-1)q}{2}$，利润 $(q^D)^2$），惩罚仍为古诺 NE 利润 $\pi^c=(a/(n+1))^2$。IC：
\[
\frac{(a-nq)q}{1-\delta}\ge \left[\frac{a-(n-1)q}{2}\right]^2+\frac{\delta}{1-\delta}\left(\frac{a}{n+1}\right)^2.
\]\indent
给定 $\delta$，可维持的利润随 $q$ 递减（在 $q\le q^m$ 区间），故利润最大化等价于取\textit{最大的可行} $q$，即上式取等号的解 $q^*(\delta)$（较大根），满足
\[
(a-nq)q=\frac{1-\delta}{1}\left[\frac{(a-(n-1)q)^2}{4}+\frac{\delta}{(n+1)^2}a^2\right].
\]\indent
其端点性质：
\[
q^*(\delta_{\min})=q^m=\frac{a}{2n},\qquad
\lim_{\delta\to 0}q^*(\delta)=\frac{a}{n+1}=q^c.
\]\indent
即 $\delta$ 越小，可维持的合谋产量越接近（甚至退化为）竞争水平；当 $\delta$ 小到无法维持任何高于古诺的产出时，唯一 SPNE 即为各期古诺 NE 本身。


% ============================================================================
%  2.16  工资—就业模型的重复博弈帕累托改进
% ============================================================================
\exer{2.16}
在第 2.1.C 节分析的工资和就业模型中，逆向归纳解对社会来讲并不是有效率的。并且在实践中，企业和工会今天就一个 3 年期合同的条款进行谈判，并将在 3 年以后再就下一期的合同条款进行谈判，如此等等。从而，两者之间的关系用重复博弈来表示可能更为准确，参见（埃斯皮诺萨和里，1989）。这个问题可导出在无限重复博弈中子博弈精炼纳什均衡帕累托优于一次性博弈中逆向归纳解的条件。分别用 $U^*$ 和 $\pi^*$ 表示一次性博弈逆向归纳解中工会效用和企业利润。考虑另外的一个工资——就业组合 $(w,L)$ 以及与之相应的效用——利润组合 $(U,\pi)$，假设双方的贴现因子都是 $\delta$（每 3 年为一个阶段），请推导 $(w,L)$ 应满足的条件，使得（1）$(U,\pi)$ 帕累托优于 $(U^*,\pi^*)$ 且（2）$(U,\pi)$ 为无限重复博弈的一个子博弈精炼纳什均衡的结果，其中只要发生任何偏离就将永远转向 $(U^*,\pi^*)$。

\sol 记 $d=a-w_0>0$。在一次性阶段博弈中，给定工资 $w$，企业最优雇佣为
\[
  L^D(w)=\frac{a-w}{2},\qquad \pi^D(w)=\frac{(a-w)^2}{4}
\]\indent
（在本题相关的 $w<a$ 范围内）。一次性逆向归纳结果为
\[
  w^*=\frac{a+w_0}{2},\quad L^*=\frac{d}{4},\quad
  U^*=\frac{d^2}{8},\quad \pi^*=\frac{d^2}{16}.
\]

现在考察一个“任何一方偏离后，下一阶段起永久转回 $(w^*,L^*)$”的触发战略。路径上约定工资—就业组合 $(w,L)$，对应当期收益
\[
  U=(w-w_0)L,\qquad \pi=(a-w-L)L.
\]\indent
若企业在路径上偏离就业水平，当前工资 $w$ 已固定，其当期最优偏离收益为 $\pi^D(w)$；此后每期得到 $\pi^*$。所以企业不偏离的条件是
\[
  \frac{\pi}{1-\delta}\ge \pi^D(w)+\frac{\delta\pi^*}{1-\delta}
  \quad\Longleftrightarrow\quad
  \pi\ge (1-\delta)\pi^D(w)+\delta\pi^*.
  \tag{1}
\]\indent
若工会单方面改报工资 $w'$, 企业在该期对 $w'$ 作最优反应，工会当期至多获得
\[
  U^D(w')=\frac{(w'-w_0)(a-w')}{2}\le U^*.
\]\indent
偏离后的未来收益也为 $U^*$，故工会不偏离的条件为
\[
  \frac{U}{1-\delta}\ge U^D(w')+\frac{\delta U^*}{1-\delta}\quad\forall w'
  \quad\Longleftrightarrow\quad U\ge U^*.
  \tag{2}
\]\indent
此外，帕累托改进要求 $U\ge U^*$、$\pi\ge\pi^*$，并至少一方严格提高。因此，这一触发战略下的可持续帕累托改进条件是
\[
  \boxed{
  (w-w_0)L\ge\frac{d^2}{8},\qquad
  (a-w-L)L\ge
  \max\!\left\{\frac{d^2}{16},\ (1-\delta)\frac{(a-w)^2}{4}+\delta\frac{d^2}{16}\right\},
  }
  \tag{3}
\]\indent
其中还须满足可行性条件 $L\ge0$（以及模型给定的工资/就业约束）。至少一个帕累托不等式须为严格不等式，才能称为严格帕累托改进。

例如取 $\delta=0.9$、$w=w_0+\frac25d$、$L=\frac13d$，则
\[
  U=\frac{2}{15}d^2>\frac18d^2=U^*,\qquad
  \pi=\frac{4}{45}d^2>\frac1{16}d^2=\pi^*.
\]\indent
企业偏离收益为 $\pi^D=\frac{9}{100}d^2$，其激励相容右侧为
\[
  (1-\delta)\pi^D+\delta\pi^*
  =\left(\frac9{1000}+\frac9{160}\right)d^2
  =\frac{261}{4000}d^2<\frac4{45}d^2=\pi.
\]\indent
工会的约束也满足。因此这确实是一个双方严格受益、且能由题设触发战略维持的均衡路径。

原先“可持续结果必须满足 $L=L^D(w)$”的推断不成立：重复博弈允许路径上的就业不同于企业当期静态最优反应，只要企业的当期偏离收益没有超过由未来惩罚带来的损失，即满足 (1) 即可。故本题存在严格帕累托改进；结论并非“只有原逆向归纳解”。

\exer{2.17}
考虑如下的一家企业和一系列工人之间的无限期博弈，每名工人只出现在一个阶段。在每一阶段工人选择努力工作，这时努力成本为 $c$，生产出产出为 $y$；或不做任何努力，也没有任何产出，不发生任何成本。如果生产出了产出，企业拥有成果但可以通过支付工资的方式与工人分享，具体方法将在后面提到。假定在一个阶段的开始工人可选择另外的就业机会，（减去努力成本的）净价值为 $0$，并且工人不能被强迫接受一个小于 $0$ 的工资。同时假定 $y>c$，从而努力工作是高效率的。在每一阶段内，博弈进行的时间顺序如下：首先工人选择一个努力水平，然后企业和工人都可以观察到产出的情况，最后企业选择支付给工人的工资水平。假定事先不能达成任何有约束力的工资协议：企业选择的工资是完全不受限制的。那么，在单一阶段的博弈中，子博弈精炼意味着无论工人的产出如何，企业给出的工资都为 $0$，于是工人也不会进行任何努力。现在考虑无限期时的情况，我们已提到每个工人只出现在一个阶段。不过我们可以假定在阶段 $t$ 的开始，博弈前面 $t-1$ 个阶段的进行过程都可被将在 $t$ 阶段工作的工人观测到（可以想象到，这些知识在工人之间代代相传）。假设企业未来收益每一阶段的贴现因子为 $\delta$。写出无限期博弈中一个子博弈精炼均衡下企业和每个工人的战略，使得对足够大的贴现因子 $\delta$，在均衡情况下，每个工人都努力工作，生产出产出 $y$，并给出上述均衡存在的充分必要条件。

\sol
记企业愿意支付的高工资为 $w>0$（待定）。构造如下战略组合 $\sigma$：

\textbf{企业战略}：在第 1 阶段支付工资 $w$；在第 $t>1$ 阶段，若此前每一阶段中，凡产出为 $y$ 时企业都支付了 $w$（即“按劳付酬”的声誉 intact），则第 $t$ 阶段仍支付 $w$；一旦曾在某阶段工人产出为 $y$ 而企业未支付 $w$，则在所有后续阶段永久支付 $0$。

\textbf{工人战略}：在第 $t$ 阶段，若企业截至 $t-1$ 阶段始终保持“支付 $w$”的声誉，则该工人选择努力（产出 $y$）；若声誉已破灭（企业曾有“不付 $w$”的记录），则选择不努力（产出 $0$）。

\textbf{工人（第 $t$ 阶段）的激励}。若声誉 intact，工人知道：自己努力则企业（为维持声誉）会支付 $w$，净收益 $w-c$；不努力则得 $0$。故努力占优当且仅当
\[
w\ge c.\tag{1}
\]

\textbf{企业的激励}。若声誉 intact，企业在第 $t$ 阶段面对努力的工人，有两种选择：
\begin{enumerate}[leftmargin=*]
  \item 守约支付 $w$：当期利润 $y-w$，此后每期继续以 $(y-w)$ 重复（声誉保持），现值
        $\displaystyle V_{\text{守}}=\frac{y-w}{1-\delta}$。
  \item 违约支付 $0$：当期利润 $y$，但声誉破灭，此后永久为 $0$，现值
        $\displaystyle V_{\text{违}}=y$。
\end{enumerate}
守约优于违约当且仅当
\[
\frac{y-w}{1-\delta}\ge y
\;\Longrightarrow\;
y-w\ge y(1-\delta)
\;\Longrightarrow\;
w\le y\delta.\tag{2}
\]

\textbf{存在的充分必要条件}。由 (1)、(2) 需 $c\le w\le y\delta$，存在这样的 $w$ 当且仅当
\[
\boxed{c\le y\delta\quad\Longleftrightarrow\quad \delta\ge \frac{c}{y}}.
\]\indent
由于 $y>c$，该阈值小于 $1$。

\textbf{结论}。当 $\delta\ge c/y$ 时，存在纯战略 SPNE 使每个工人都努力工作、企业每期支付 $w\in[c,\,y\delta]$；当 $\delta<c/y$ 时，任何 SPNE 都退化为“企业支付 $0$、工人不努力”的低效率结果。$\blacksquare$

% ============================================================================
%  2.18  一般动态博弈中的概念
% ============================================================================
\exer{2.18}
对一个任意的博弈，解释其战略、信息集以及子博弈的概念。

\sol
\begin{enumerate}[leftmargin=*,label=(\roman*)]
  \item \textbf{战略（Strategy）}。参与者 $i$ 的一个（纯）战略，是其针对\textit{每一个自己可能到达的信息集}规定一个行动（或行动的概率分布，若考虑混合/行为战略）的完整规则。关键点在于：战略必须是“事前”就完备的——它要说明在博弈树上\textit{每一个}该参与者可能轮到行动的节点（信息集）上该怎么做，而不仅仅是在均衡路径上的节点。一个战略组合决定了博弈树的唯一一条路径及相应收益。
  \item \textbf{信息集（Information set）}。信息集是博弈树上归属于同一参与者、且该参与者在轮到行动时\textit{无法区分}的一组节点的并集。同一信息集内的节点必须满足：(a) 都属于同一参与者；(b) 该参与者在每个节点上可选择的行动集合相同。信息集体现了“同时行动”或“不可观测他人先前行动”所带来的不完美信息。若某信息集只含一个节点，称该信息集是“单点”的（完美信息）；否则为不完美信息。
  \item \textbf{子博弈（Subgame）}。原博弈的一个子博弈，是从博弈树的某个\textit{单点信息集}（即单个节点）开始、并包含该节点之后整棵后续子树（含所有后代节点与信息集）的一个子博弈。要求起点必须是单点信息集，是因为子博弈的参与者必须从一开始就知道“我们确实进入了这个子博弈”。初始节点本身当然也是单点信息集，故整个博弈是最大的子博弈。
\end{enumerate}
\cmt{子博弈精炼纳什均衡要求战略组合在每个子博弈（而不仅在整个博弈）中都是纳什均衡——这等价于排除了依靠“不可置信的威胁”维持的均衡。}


% ============================================================================
%  2.19  三阶段鲁宾斯坦讨价还价的 SPNE
% ============================================================================
\exer{2.19}
在第 2.1.D 节分析的三阶段鲁宾斯坦讨价还价模型中，我们解出了逆向归纳解。它的子博弈精炼纳什均衡是什么？

\sol
三阶段模型（两参与者贴现因子均为 $\delta$）：
\begin{enumerate}[leftmargin=*]
  \item $t=1$：参与者 1 出价 $(x_1,1-x_1)$；
  \item 若拒绝，$t=2$：参与者 2 出价 $(x_2,1-x_2)$；
  \item 若再拒绝，$t=3$：参与者 1 出价 $(x_3,1-x_3)$；
  \item 若再拒绝，博弈结束，双方收益 $(0,0)$。
\end{enumerate}
由于博弈有限且每个子博弈的逆向归纳唯一，逆向归纳解即 SPNE。从后往前：
\begin{enumerate}[leftmargin=*,label=(\roman*)]
  \item \textbf{$t=3$（参与者 1 出价）}：拒绝则 $(0,0)$，故参与者 1 提出 $(1,0)$，参与者 2 接受。该子博弈结果为 $(1,0)$。
  \item \textbf{$t=2$（参与者 2 出价）}：若参与者 2 的报价被拒，进入 $t=3$，参与者 1 将得 $1$（贴现到 $t=2$ 为 $\delta\cdot 1$）。因此参与者 2 只需给参与者 1 份额 $\delta$ 即可使其接受，自己保留 $1-\delta$。结果为 $(\delta,\,1-\delta)$。
  \item \textbf{$t=1$（参与者 1 出价）}：若被拒，进入 $t=2$，参与者 2 将得 $1-\delta$（贴现到 $t=1$ 为 $\delta(1-\delta)$）。故参与者 1 给参与者 2 份额 $\delta(1-\delta)$，自己保留
  \[
  x_1^*=1-\delta(1-\delta)=1-\delta+\delta^2.
  \]\indent
  结果为 $(1-\delta+\delta^2,\ \delta-\delta^2)$。
\end{enumerate}

\textbf{SPNE 战略组合}（以“坚持某要价、达到某最低接受门槛即接受”的形式写出）：
\begin{align*}
\sigma_1:\ &t=1\text{ 出价 }(1-\delta+\delta^2,\ \delta-\delta^2);\quad
            t=3\text{ 出价 }(1,0);\\
           &\text{在 }t=1,t=3\text{ 作为回应方时不出现，故无回应规则}\\[-2pt]
           &\text{（或：接受任何使其份额}\ge\delta\text{ 的报价）。}\\
\sigma_2:\ &t=2\text{ 出价 }(\delta,\ 1-\delta);\\
           &\text{在 }t=1\text{ 接受当且仅当其份额}\ge \delta(1-\delta);\quad
            \text{在 }t=3\text{ 接受当且仅当其份额}\ge 0.
\end{align*}\indent
该战略组合即为唯一的子博弈精炼纳什均衡（标准 tie-breaking 下），逆向归纳解 $(1-\delta+\delta^2,\ \delta-\delta^2)$ 正是其在第 1 期的产出。\cmt{当 $\delta\to1$ 时 $x_1^*\to1$；当 $\delta\to0$ 时 $x_1^*\to1$，但中间 $\delta=1/2$ 时 $x_1^*=3/4$，体现先动者优势随时间贴现而波动。}

% ============================================================================
%  2.20  无限期鲁宾斯坦讨价还价的 SPNE 证明
% ============================================================================
\exer{2.20}
在无限期鲁宾斯坦讨价还价模型中考虑下面的战略（注意习惯性的表示方法，开价 $(s,1-s)$ 意味着参与者 1 将会得到 $s$，参与者 2 将得到 $1-s$，而不论是哪一方提出的条件）：令 $s^*=1/(1+\delta)$，参与者 1 坚持开价 $(s^*,1-s^*)$，并且只有当 $s\ge \delta\cdot s^*$ 时，才接受对方开价 $(s,1-s)$；参与者 2 坚持开价 $(1-s^*,s^*)$，并且只有当 $1-s\ge \delta\cdot s^*$ 时，才接受对方开价 $(s,1-s)$。证明这两个战略是一个纳什均衡，并证明这一均衡是子博弈精炼的。

\sol
记 $\underline{s}=\delta s^*=\dfrac{\delta}{1+\delta}$，则参与者 1 接受条件为 $s\ge \underline{s}$，参与者 2 接受条件为 $1-s\ge \underline{s}\iff s\le 1-\underline{s}=s^*$。

\textbf{（一）证明是纳什均衡}。设参与者 2 遵循 $\sigma_2$。考察参与者 1 的各类可能偏离：
\begin{enumerate}[leftmargin=*,label=(\alph*)]
  \item 当轮到参与者 1 出价时，若他偏离 $\sigma_1$ 提出别的 $(s,1-s)$：
        \begin{enumerate}[leftmargin=*]
          \item 若 $s<s^*$（给参与者 2 更多），参与者 2 仍接受（因 $1-s>s^*$），参与者 1 所得 $s<s^*$，变差；
          \item 若 $s>s^*$（给参与者 2 更少，即 $1-s<\underline{s}$），参与者 2 拒绝；轮到参与者 2 出价 $(1-s^*,s^*)$，参与者 1 得 $1-s^*=\underline{s}<s^*$，变差。
        \end{enumerate}
        故参与者 1 在出价时不会偏离。
  \item 当轮到参与者 1 回应参与者 2 的报价 $(s,1-s)$ 时，其策略要求接受当且仅当 $s\ge \underline{s}$。若他拒绝一个本应接受的报价（$s\ge\underline{s}$），则下一期轮到他出价得到 $s^*$，而 $s^*\ge s$ 且拒绝损失一期贴现，故拒绝后所得 $\delta s^*=\underline{s}\le s$，不会更好；若他接受一个本应拒绝的报价（$s<\underline{s}$），则直接得 $s<\underline{s}=\delta s^*$，比拒绝后出价得 $s^*$ 贴现 $\delta s^*$ 更差。故回应规则最优。
\end{enumerate}
由对称性，参与者 2 在 $\sigma_1$ 下亦无有利可图的偏离。因此 $\sigma=(\sigma_1,\sigma_2)$ 是纳什均衡。

\textbf{（二）证明是子博弈精炼的}。对任意历史（子博弈），$\sigma$ 规定的行动只取决于“当前由谁出价”以及“当前报价”，与历史无关；且上述一步偏离分析在\textit{每一个}出价节点与回应节点上都成立（因为各子博弈的“未来贴现价值”均为同样的 $s^*$ 或 $1-s^*$）。具体看四类子博弈：
\begin{enumerate}[leftmargin=*]
  \item 参与者 1 出价节点：提出 $(s^*,1-s^*)$ 最优（如上），且参与者 2 的接受规则使该报价被接受；
  \item 参与者 2 出价节点：提出 $(1-s^*,s^*)$ 最优（对称论证）；
  \item 参与者 1 回应节点：接受门槛 $s\ge\underline{s}$ 恰使其对“拒绝后下一期得 $s^*$”无偏离激励；
  \item 参与者 2 回应节点：接受门槛 $1-s\ge\underline{s}$ 同理最优。
\end{enumerate}
故 $\sigma$ 在每个子博弈都构成纳什均衡，因而是子博弈精炼纳什均衡。$\blacksquare$

% ============================================================================
%  2.21  手雷博弈（扩展式、标准式、NE / 逆向归纳 / SPNE）
% ============================================================================
\exer{2.21}
给出第 2.1 节描述的手雷博弈的扩展式表述及标准式表述。并分别写出其纯战略纳什均衡、逆向归纳解和子博弈精炼纳什均衡。

\sol
\textbf{模型（三个阶段，参与者 1 行动两次）}。
\begin{enumerate}[leftmargin=*,label=(\arabic*)]
  \item 参与者 1 选择 $L$ 或 $R$。若选 $L$，博弈结束，收益为 $(2,0)$。
  \item 若参与者 1 选 $R$，参与者 2 观测到后选择 $L'$ 或 $R'$。若选 $L'$，博弈结束，两人的收益均为 $1$。
  \item 若参与者 2 选 $R'$，参与者 1（记得自己此前选了 $R$）再选择 $L''$ 或 $R''$，无论选哪个博弈都结束：选 $L''$ 得 $(3,0)$，选 $R''$ 得 $(0,2)$。
\end{enumerate}

\textbf{扩展式}。三阶段完美信息博弈树：

\myfig{%
\begin{tikzpicture}[scale=0.92, every node/.style={font=\small}]
  \node[pnode] (p1) at (0,0) {};
  \node[above=2pt of p1, nodlab] {$1$};
  \node[tnode] (t1) at (-2.6,-2.6) {};
  \node[below=3pt of t1, pay] {$2,\ 0$};
  \node[pnode] (p2) at (2.6,-2.6) {};
  \node[right=4pt of p2, nodlab] {$2$};
  \node[tnode] (t2) at (0,-5.0) {};
  \node[below=3pt of t2, pay] {$1,\ 1$};
  \node[pnode] (p3) at (5.4,-5.0) {};
  \node[right=4pt of p3, nodlab] {$1$};
  \node[tnode] (t3) at (2.8,-7.4) {};
  \node[below=3pt of t3, pay] {$3,\ 0$};
  \node[tnode] (t4) at (8.0,-7.4) {};
  \node[below=3pt of t4, pay] {$0,\ 2$};
  \draw[branch] (p1) -- (t1) node[pos=0.5, elab, left=1pt] {$L$};
  \draw[branch] (p1) -- (p2) node[pos=0.5, elab, right=1pt] {$R$};
  \draw[branch] (p2) -- (t2) node[pos=0.5, elab, left=1pt] {$L'$};
  \draw[branch] (p2) -- (p3) node[pos=0.5, elab, right=1pt] {$R'$};
  \draw[branch] (p3) -- (t3) node[pos=0.5, elab, left=1pt] {$L''$};
  \draw[branch] (p3) -- (t4) node[pos=0.5, elab, right=1pt] {$R''$};
\end{tikzpicture}%
}{手雷博弈的扩展式}

\textbf{标准式}。参与者 1 有两个决策点，故有 $2\times2=4$ 个纯战略，记为 $(a_1,a_3)\in\{L,R\}\times\{L'',R''\}$；参与者 2 只有一个决策点（位于参与者 1 选 $R$ 之后），有 $2$ 个纯战略 $L',R'$。收益表为

\begin{paytab}
\begin{tabular}{c|cc}
\pthead{$1\backslash 2$} & $L'$ & $R'$\\\mapline
$(L,L'')$  & $2,0$ & $2,0$\\\mapline
$(L,R'')$  & $2,0$ & $2,0$\\\mapline
$(R,L'')$  & $1,1$ & $3,0$\\\mapline
$(R,R'')$  & $1,1$ & $0,2$\\\mapline
\end{tabular}
\end{paytab}

\textbf{纯战略纳什均衡}。对参与者 2：给定参与者 1 的任一战略，选 $L'$ 的收益为 $0$（若 1 选 $L$）、$1$（若 1 选 $R$）；选 $R'$ 的收益为 $0$（若 1 选 $L$）、$3$ 的对手方收益下 2 得 $0$（若 1 选 $(R,L'')$）或 $2$（若 1 选 $(R,R'')$）。逐列考察参与者 1 的最优反应：
\begin{enumerate}[leftmargin=*]
  \item 在 $L'$ 列，参与者 1 的收益依次为 $2,2,1,1$，最优为 $(L,L'')$ 或 $(L,R'')$；
  \item 在 $R'$ 列，收益依次为 $2,2,3,0$，最优为 $(R,L'')$。
\end{enumerate}
对参与者 2：给定 $(L,L'')$ 或 $(L,R'')$（博弈在第一阶段即结束），$L'$ 与 $R'$ 均给出 $0$，两者都是最优反应；给定 $(R,L'')$，$L'$ 给 $1$ 而 $R'$ 给 $0$，故最优为 $L'$；给定 $(R,R'')$，$L'$ 给 $1$ 而 $R'$ 给 $2$，故最优为 $R'$。

两两配对可得纯战略纳什均衡恰为
\[
\boxed{\bigl((L,L''),\,L'\bigr)\quad\text{与}\quad\bigl((L,R''),\,L'\bigr),}
\]\indent
两者产出均为 $(2,0)$。

\textbf{逆向归纳解}。从第三阶段倒推：参与者 1 面临 $L''$（得 $3$）与 $R''$（得 $0$），故选 $L''$。预见此点，阶段 2 参与者 2 选 $L'$ 得 $1$，选 $R'$ 则最终得 $0$，故选 $L'$。预见此点，阶段 1 参与者 1 选 $L$ 得 $2$，选 $R$ 则最终得 $1$，故选 $L$。故逆向归纳解为
\[
(L,\ L',\ L''),\qquad\text{产出 }(2,0).
\]

\textbf{子博弈精炼纳什均衡}。SPNE 要求在每个子博弈上都最优。第三阶段的子博弈迫使参与者 1 选 $L''$，这排除了含 $R''$ 的战略；第二阶段的子博弈迫使参与者 2 选 $L'$；第一阶段参与者 1 选 $L$。故唯一 SPNE 为
\[
\boxed{\bigl((L,L''),\,L'\bigr),\qquad\text{产出 }(2,0).}
\]\indent
纯战略 NE $\bigl((L,R''),L'\bigr)$ 不是子博弈精炼的：它要求参与者 1 在第三阶段承诺选择 $R''$（得 $0$），而一旦真的到达该节点他就会改选 $L''$（得 $3$）——这正是一个不可置信的威胁。

\cmt{本博弈是“承诺的价值”的经典例子：参与者 1 在第一阶段选择 $L$（得 $2$）看似放弃了第三阶段可能得到的 $3$，但正是这一“先行承诺”使参与者 2 在第二阶段愿意接受 $1$ 而选择 $L'$，从而让 1 保住了 $2$。反之，若 1 在第一阶段选 $R$，则 2 会选 $R'$，1 最终只能得 $0$——逆向归纳之所以要“向前看”，正是为了避免这一后果。}

% ============================================================================
%  2.22  银行挤提博弈（扩展式、标准式、纯战略 SPNE）
% ============================================================================
\exer{2.22}
给出第 2.2.B 节讨论的银行挤提博弈的扩展式表述及标准式表述。其纯战略子博弈精炼纳什均衡是什么？

\sol
\textbf{模型}。两位投资者各存入 $D$；银行把存款投入长期项目。若被迫在项目到期前变现，共可收回 $2r$，其中 $D>r>D/2$；若允许项目到期，则共可取得 $2R$，其中 $R>D$。不设贴现。

第 1 期（到期前）两人\emph{同时}决定提款（$W$）或不提（$N$）：
\begin{enumerate}[leftmargin=*]
  \item 两人\emph{都}提：项目被变现，各得 $r$，博弈结束；
  \item 只有 $i$ 提：他得 $D$，另一人得 $2r-D$，博弈结束；
  \item 两人都\emph{不}提：项目到期，进入第 2 期。
\end{enumerate}
第 2 期（到期后）两人同时决定提或不提：
\begin{enumerate}[leftmargin=*]
  \item 两人都提：各得 $R$；
  \item 只有 $i$ 提：他得 $2R-D$，另一人得 $D$；
  \item 两人都不提：银行各返还 $R$。
\end{enumerate}
两个阶段的标准式分别为

\begin{paytab}
\begin{tabular}{c|cc}
\pthead{日期 1} & $W$ & $N$\\\mapline
$W$ & $r,\,r$ & $D,\,2r-D$\\\mapline
$N$ & $2r-D,\,D$ & \textnormal{（进入第 2 期）}\\\mapline
\end{tabular}
\end{paytab}
\begin{paytab}
\begin{tabular}{c|cc}
\pthead{日期 2} & $W$ & $N$\\\mapline
$W$ & $R,\,R$ & $2R-D,\,D$\\\mapline
$N$ & $D,\,2R-D$ & $R,\,R$\\\mapline
\end{tabular}
\end{paytab}

\textbf{扩展式}。每一期的两个决策节点属于同一信息集（同时行动，故参与者 2 观测不到参与者 1 当期的选择）：

\myfig{%
\begin{tikzpicture}[scale=0.9, every node/.style={font=\small}]
  \node[pnode] (a) at (0,0) {};
  \node[above=2pt of a, nodlab] {$1$};
  \node[pnode] (b) at (3.2,-2.8) {};
  \node[right=4pt of b, nodlab] {$2$};
  \node[pnode] (c) at (-3.2,-2.8) {};
  \node[left=4pt of c, nodlab] {$2$};
  \node[tnode] (ww) at (4.8,-4.8) {}; \node[below=3pt of ww, pay] {$(r,r)$};
  \node[tnode] (wn) at (1.7,-4.8) {}; \node[below=3pt of wn, pay] {$(D,\,2r-D)$};
  \node[tnode] (nw) at (-1.7,-4.8) {}; \node[below=3pt of nw, pay] {$(2r-D,\,D)$};
  \node[tnode] (nn) at (-4.8,-4.8) {}; \node[below=3pt of nn, pay] {进入第 2 期};
  \draw[branch] (a) -- (b) node[pos=0.5, elab, above right=-1pt] {$W$};
  \draw[branch] (a) -- (c) node[pos=0.5, elab, above left=-1pt] {$N$};
  \draw[branch] (b) -- (ww) node[pos=0.5, elab, right=1pt] {$W$};
  \draw[branch] (b) -- (wn) node[pos=0.5, elab, left=1pt] {$N$};
  \draw[branch] (c) -- (nw) node[pos=0.5, elab, right=1pt] {$W$};
  \draw[branch] (c) -- (nn) node[pos=0.5, elab, left=1pt] {$N$};
  \draw[infoset] (c) -- (b);
  % ---- 第 2 期子树 ----
  \node[pnode] (a2) at (-4.8,-7.6) {}; \node[above=2pt of a2, nodlab] {$1$};
  \node[pnode] (b2) at (-2.4,-9.4) {}; \node[right=3pt of b2, nodlab] {$2$};
  \node[pnode] (c2) at (-7.2,-9.4) {}; \node[left=3pt of c2, nodlab] {$2$};
  \node[tnode] (w2) at (-1.3,-11.2) {}; \node[below=3pt of w2, pay] {$(R,R)$};
  \node[tnode] (n2) at (-3.5,-11.2) {}; \node[below=3pt of n2, pay] {$(2R-D,\,D)$};
  \node[tnode] (v2) at (-6.1,-11.2) {}; \node[below=3pt of v2, pay] {$(D,\,2R-D)$};
  \node[tnode] (m2) at (-8.3,-11.2) {}; \node[below=3pt of m2, pay] {$(R,R)$};
  \draw[branch] (a2) -- (b2) node[pos=0.5, elab, above right=-1pt] {$W$};
  \draw[branch] (a2) -- (c2) node[pos=0.5, elab, above left=-1pt] {$N$};
  \draw[branch] (b2) -- (w2) node[pos=0.5, elab, right=1pt] {$W$};
  \draw[branch] (b2) -- (n2) node[pos=0.5, elab, left=1pt] {$N$};
  \draw[branch] (c2) -- (v2) node[pos=0.5, elab, right=1pt] {$W$};
  \draw[branch] (c2) -- (m2) node[pos=0.5, elab, left=1pt] {$N$};
  \draw[infoset] (c2) -- (b2);
  \draw[contarc] (-4.8,-5.5) -- (-4.8,-7.2);
\end{tikzpicture}%
}{银行挤提博弈的扩展式（第 1 期与第 2 期）}

\textbf{标准式}。每位投资者在第 1 期有一个行动、在第 2 期有一个行动，故每个投资者有 $2\times2=4$ 个纯战略。以 $W_1W_2$ 表示“第 1 期提、第 2 期提”，$W_1N_2$ 表示“第 1 期提、第 2 期不提”，$N_1W_2$ 与 $N_1N_2$ 类似。完整 $4\times4$ 收益矩阵（行=投资者 1，列=投资者 2）为

\begin{paytab}
\begin{tabular}{c|cccc}
\pthead{$1\backslash2$} & $W_1W_2$ & $W_1N_2$ & $N_1W_2$ & $N_1N_2$\\\mapline
$W_1W_2$ & $(r,r)$ & $(r,r)$ & $(D,\,2r-D)$ & $(D,\,2r-D)$\\\mapline
$W_1N_2$ & $(r,r)$ & $(r,r)$ & $(D,\,2r-D)$ & $(D,\,2r-D)$\\\mapline
$N_1W_2$ & $(2r-D,\,D)$ & $(2r-D,\,D)$ & $(R,R)$ & $(2R-D,\,D)$\\\mapline
$N_1N_2$ & $(2r-D,\,D)$ & $(2r-D,\,D)$ & $(D,\,2R-D)$ & $(R,R)$\\\mapline
\end{tabular}
\end{paytab}
（前两行中第 2 期行动不起作用，因为两人在第 1 期就提款、项目被变现，博弈已结束。）

\textbf{逆向归纳：第 2 期子博弈}。由 $R>D$ 得 $2R-D>R>D$。若对方提款，我提款得 $R$、不提得 $D$，故提款更优；若对方不提，我提款得 $2R-D$、不提得 $R$，因 $2R-D>R$ 仍应提款。故“提款”严格占优，第 2 期的唯一纳什均衡是
\[
(W,W)\ \text{产出}\ (R,R).
\]\indent
任何子博弈精炼均衡在第 2 期都必须规定双方提款。

\textbf{逆向归纳：回到第 1 期}。把第 2 期的解代入，“两人都不提”这一格变为 $(R,R)$，第 1 期Reduced 收益表为
\[
\begin{array}{c|cc}
 & W & N\\\hline
W & (r,r) & (D,\,2r-D)\\
N & (2r-D,\,D) & (R,R)
\end{array}
\]\indent
\begin{enumerate}[leftmargin=*]
  \item 对方提款时：我提款得 $r$，不提得 $2r-D$。因 $D>r$，有 $r-(2r-D)=D-r>0$，故提款是唯一最优反应。于是 $(W,W)$ 是纳什均衡，产出 $(r,r)$。
  \item 对方不提款时：我提款得 $D$，不提得 $R$。因 $R>D$，故不提是唯一最优反应。于是 $(N,N)$ 是纳什均衡，产出 $(R,R)$。
\end{enumerate}
第 1 期恰有两个纯战略纳什均衡，且各自在第 2 期都接上唯一子博弈解，故两者都是子博弈精炼的。纯战略 SPNE 恰为两个（把战略写成“第 1 期行动, 第 2 期行动”）：
\[
\boxed{\bigl((W,W),(W,W)\bigr)\ \text{产出 }(r,r)
\qquad\text{与}\qquad
\bigl((N,W),(N,W)\bigr)\ \text{产出 }(R,R).}
\]\indent
其余战略组合都不是纳什均衡：若只有一方在第 1 期提款，\emph{提款方}可以改为不提而把结果从 $D$ 提升到 $R$（对方不提时），故其最优反应必然是不提，与“只有一方提”的历史矛盾。

\cmt{第一个均衡是“坏”均衡（无效率的挤提，各得 $r<R$），第二个是“好”均衡（持有至到期，各得 $R$）。这正是银行挤提多重均衡的本质：完全信息、完美理性之下仍可能出现自我实现的信心危机——两人只要都相信对方会挤提，挤提就是理性选择。注意本题的关键在于 $D>r$ 与 $R>D$ 这两个方向相反的激励：$D>r$ 使“对方提款时我也提款”成为最优反应（挤提均衡），$R>D$ 使“对方不提款时我也不提”成为最优反应（不挤提均衡），二者共同保证了上述两个均衡各自稳定。}

% ============================================================================
%  2.23  买方专用投资、卖方不可观测时的交易博弈
% ============================================================================
\exer{2.23}
一个卖方和一个买方打算进行交易。在他们交易之前，买方可以作一项投资，从而提高标的物对他的价值。这项投资不能被卖方观察到，从而也不会影响标的物对卖方的价值，后者我们标准化为 $0$。购买方对标的的初始价值为 $v>0$；一项投资 $I$ 使得购买方价值变为 $v+1$，但相应增加了成本 $I^2$。博弈进行的时间顺序如下：首先，购买方选择投资水平 $I$，发生成本 $I^2$；第二，卖方不能观测到 $I$，但开出卖价为 $p$；第三，买方或者接受，或者拒绝卖方的开价：如果买方接受，则买方的收益为 $v+1-p-I^2$，卖方的收益为 $p$；如果拒绝，则双方的收益分别为 $-I^2$ 和 $0$。证明这一博弈不存在纯战略子博弈精炼纳什均衡。解出博弈的混合战略纳什均衡，其中买方的混合战略中，出现概率为正的只有两种投资水平，并且卖方的混合战略中，出现概率为正的只有两个价格水平。

\sol
\textbf{第一步：买方的接受\/拒绝决策。} 给定卖方开出的价格 $p$ 与既定的投资 $I$，买方比较"接受"与"拒绝"的收益：
\[
\text{接受：}v+1-p-I^2,\qquad \text{拒绝：}-I^2.
\]\indent
接受不劣于拒绝当且仅当
\[
v+1-p-I^2\ge -I^2\;\Longleftrightarrow\; p\le v+1.
\]\indent
关键的观察是：$I^2$ 在比较中被消掉了。因此买方的接受\/拒绝决策\textbf{只取决于价格 $p$ 是否满足 $p\le v+1$}，与投资水平 $I$ 无关。

\textbf{第二步：为何不存在纯战略子博弈精炼纳什均衡。} 固定任意卖方战略 $p$。由于卖方不能观测 $I$，其决策节点之后买方仍可自由选择 $I$。考察买方在\emph{接受}场合（$p\le v+1$）的收益
\[
\pi_B(I)=v+1-p-I^2,\qquad I>0,
\]\indent
它关于 $I$ 严格递减，其\emph{上确界}在 $I\to 0^+$ 处取得：
\[
\sup_{I>0}\pi_B(I)=v+1-p,
\]\indent
但该上确界\emph{在纯战略中不可达}——任何 $I>0$ 实际只给 $v+1-p-I^2<v+1-p$，而 $I=0$ 虽可达却不带来 $+1$ 的增值（价值仍为 $v$）。同样地，在\emph{拒绝}场合收益 $-I^2$ 也在 $I\to0^+$ 处取上确界 $0$ 而不可达。

因此，无论卖方定什么价，买方在投资维度上的收益都只有上确界而无最大值，\textbf{不存在纯战略最佳反应}。既然子博弈精炼纳什均衡要求在每一个子博弈后（包括卖方决策之后的子博弈）都取到最佳反应，而此处最佳反应不存在，故\emph{本博弈不存在纯战略子博弈精炼纳什均衡}。

\cmt{本问的关键机制：题目设定中投资的回报是"离散的整块加成"（价值从 $v$ 跳到 $v+1$），而成本 $I^2$ 是连续函数。这一"加成 + 连续成本"的组合使买方收益的上确界落在 $I\to0^+$ 而非任何可取的纯战略上，正是这一"不可达的上确界"导致纯策略均衡不存在。}

\textbf{第三步：混合战略纳什均衡的构造。} 由于纯战略不可行，均衡只能建立在混合策略上：买方以一定概率在两个投资水平间随机化，卖方以一定概率在两个价格水平间随机化，使双方在各自的支持集内无差异。

设卖方以概率 $\theta$ 开高价 $p_H$、以概率 $1-\theta$ 开低价 $p_L$，其中
\[
v<p_H\le v+1,\qquad p_L\le v.
\]\indent
这样安排的激励含义是：低价 $p_L$ 对两类买方都有吸引力（都愿接受），高价 $p_H$ 只有"高增值"买方（即取得 $+1$ 加成的投资）才接受，未投资的买方因 $p_H>v$ 而拒绝。记买方选择"投资"（取 $I_H>0$，价值 $v+1$）的概率为 $\mu$。

\begin{enumerate}[leftmargin=2em]
  \item \textbf{卖方无差异条件。} 卖方开低价 $p_L$ 时，买方\emph{无论}是否投资都会接受（因 $p_L\le v<v+1$），故卖方期望收入为 $p_L$；开高价 $p_H$ 时只有已投资方接受，买方以概率 $\mu$ 投资，故卖方期望收入为 $\mu p_H$。无差异要求
    \[
    (1-\theta)\,p_L+\theta\,(\mu p_H)=p_L.
    \]\indent
  \item \textbf{买方无差异条件。} 未投资方（$I=0$，价值 $v$）在低价下接受、高价下拒绝：
    \[
    U_{0}=(1-\theta)(v-p_L)+\theta\cdot 0=(1-\theta)(v-p_L).
    \]\indent
    已投资方（$I_H$，价值 $v+1$，成本 $I_H^2$）在两个价位下都接受：
    \[
    U_{I}=(1-\theta)\bigl(v+1-p_L-I_H^2\bigr)+\theta\bigl(v+1-p_H-I_H^2\bigr).
    \]\indent
    令二者相等：
    \begin{align*}
    (1-\theta)(v-p_L)
    &=(1-\theta)\bigl(v+1-p_L-I_H^2\bigr)+\theta\bigl(v+1-p_H-I_H^2\bigr)\\[4pt]
    &\Longleftrightarrow\;
    \theta\bigl(v+1-p_H-I_H^2\bigr)
    =(1-\theta)\,I_H^2 .
    \end{align*}\indent
\end{enumerate}

上式连同卖方的无差异条件 $p_L=(1-\theta)p_L+\theta\mu p_H$ 一起，$\bigl(\mu,\theta,I_H,p_L,p_H\bigr)$ 五个未知数由"双方无偏离"这一组条件（加上互补松弛／支持集约束）锁定，刻画出本题的混合战略纳什均衡。

\cmt{要点归纳。（i）卖方无差异条件把"开低价稳赚 $p_L$"与"开高价按 $\mu$ 概率赚 $p_H$"拉平，$\theta$ 正是这一筛选威胁的可信度；（ii）买方无差异条件表明，只有当卖方以正概率开高价（$\theta>0$）时，投资的期望增值才可能被私人成本 $I_H^2$ 抵消，从而存在使投资"值得"的概率 $\theta$；（iii）由于本问题的参数使上确界不可达，$\mu$ 与 $\theta$ 无法取到使均衡退化的角点解，故均衡是真正非退化的混合均衡，买方支持集恰为两个投资水平、卖方支持集恰为两个价格水平，与题面要求吻合。}




